% =====================================================================
%  Harald Høffding, A SHORT SURVEY OF THE HISTORY OF MODERN PHILOSOPHY.
%  English translation of: Kort Oversigt over den nyere Filosofis Historie.  Fjerde, paany omarbejdede Udgave.
%  Kjøbenhavn og Kristiania: Gyldendalske Boghandel, Nordisk Forlag, 1910.
%
%  GENERATED; DO NOT EDIT BY HAND.  Built by ebook-method/kort-oversigt/tr_parts.py build --write from
%  tr_template.tex and the parts in ebook-method/kort-oversigt/tr/ (translated 2026-10-17 from
%  texts/hoeffding/kort-oversigt/transcription.tex, following TRANSLATOR-BRIEF.md and GLOSSARY.md there).
%  Page numbers in the margin [N] are those of the Danish print.  The one footnote is marked *) as printed.
%  Title: A Short Survey of the History of Modern Philosophy (Hans, 2026-10-17); not to be confused with the longer
%  "A Brief History of Modern Philosophy" (Macmillan 1912).  Terms: ebook-method/nyere-filosofis-historie/GLOSSARY.md.
%  ODD SPOTS: the six translators listed 54 places in the Danish that looked wrong or unclear (tr/odd-*.tsv); all were
%  read at the image: 33 faults of the transcription, mended in the Danish (patches.py) and here (the first line of
%  p. 2, a lost paragraph break on p. 32, specks, doubled dashes, italic titles); 9 misprints of the print, kept and
%  marked in the Danish, translated here as the sense requires.  The rest were grammar, fact or period spelling.
% =====================================================================
\documentclass[12pt,a4paper,openany]{book}
\usepackage[utf8]{inputenc}
\usepackage[T1]{fontenc}
\usepackage{libertinus}
\usepackage{textalpha}
\usepackage[english]{babel}
\usepackage{setspace}
\usepackage[margin=1.2in]{geometry}
\usepackage{fancyhdr}
\setlength{\headheight}{28pt}
\addtolength{\topmargin}{-13.5pt}
\usepackage{hyperref}
\pagestyle{fancy}
\fancyhf{}
\fancyhead[LE,RO]{\thepage}
\fancyhead[RE]{\textit{Harald Høffding: A Short Survey}}
\fancyhead[LO]{\textit{\leftmark}}
\fancypagestyle{plain}{\fancyhf{}\fancyhead[LE,RO]{\thepage}}
\setstretch{1.3}
\usepackage{marginnote}
\usepackage{graphicx}
\DeclareUnicodeCharacter{0254}{\reflectbox{c}}  % ɔ: the old Danish id-est mark (a reversed c)
\newcommand{\starmark}[1]{\ifcase#1\or*\or**\or***\fi)}   % the print marks its one note *)
\renewcommand{\thefootnote}{\starmark{\value{footnote}}}
\newcommand{\opage}[1]{\marginnote{\footnotesize\textit{[#1]}}}
\newcommand{\apage}[1]{\marginnote{\footnotesize\textit{[#1]}}}
\newcommand{\whiteline}{\par\vspace{\baselineskip}\par}
% the printed Indhold: an essay or chapter line (title, page), a section line
\newcommand{\ientry}[2]{\par\smallskip\noindent #1\dotfill\,#2\par}
\newcommand{\isub}[2]{\par\noindent\hspace*{2em}#1\dotfill\,#2\par}
% the headings (no \addcontentsline: the print has its own Indhold, and a footnote in a heading breaks a toc entry)
\pdfstringdefDisableCommands{\renewcommand{\footnote}[2][]{}}
\newcommand{\prefacehead}[1]{\chapter*{#1}\markboth{Preface}{Preface}}
\newcommand{\sectionhead}[1]{\par\medskip\begin{center}\textbf{#1}\end{center}\smallskip}   % the print: centred, bold
\newcommand{\introhead}[1]{\begin{center}{\large #1}\\[4pt]\rule{2em}{0.4pt}\end{center}\markboth{Introduction}{Introduction}\medskip}
\newcommand{\closerule}{\par\bigskip\begin{center}\rule{4em}{0.4pt}\end{center}\par}   % the short rule that closes each part of the print
\newcommand{\chapterheadone}[2]{\bigskip\begin{center}{\large #1}\\[4pt]\rule{2em}{0.4pt}\end{center}\markboth{#2}{#2}\medskip}   % VIII: numeral and title on one line
\newcommand{\chapterhead}[2]{\bigskip\begin{center}#1\\[3pt]{\large #2}\\[4pt]\rule{2em}{0.4pt}\end{center}\markboth{#2}{#2}\medskip}   % numeral, title in capitals, short rule

\begin{document}

\frontmatter
\thispagestyle{empty}
\vspace*{5cm}
\begin{center}
{\large A SHORT SURVEY OF THE\\[4pt] HISTORY OF MODERN PHILOSOPHY}
\end{center}
\clearpage
\begin{titlepage}
\centering
\vspace*{2cm}
{HARALD HØFFDING}\par
\vspace{6pt}\rule{2em}{0.4pt}\par
\vspace{1.2cm}
{\Huge A SHORT SURVEY OF\\[6pt] THE HISTORY OF\\[6pt] MODERN PHILOSOPHY}\par
\vspace{1cm}
{\small FOURTH, NEWLY REVISED EDITION}\par
\vfill
{\small [the publisher's device]}\par
\vfill
{COPENHAGEN AND KRISTIANIA}\par
{GYLDENDALSKE BOGHANDEL}\par
{NORDISK FORLAG}\par
\vspace{4pt}\rule{2em}{0.4pt}\par\vspace{4pt}
{\footnotesize FORLAGETS BOGTRYKKERI}\par
{\footnotesize 1910}\par
\vspace{1cm}
{\small\textit{Translated from the Danish}}\par
\end{titlepage}
\clearpage

\chapter*{CONTENTS}
\markboth{Contents}{Contents}
\noindent\hfill{\small Page}\par
\ientry{Introduction. The Philosophical Problems}{1}
\ientry{I. The Origin of Modern Philosophy}{5}
\isub{1. The Discovery of Natural Man}{5}
\isub{2. The New Picture of the World}{11}
\isub{3. The New Science}{14}
\ientry{II. The Formation of Systems}{17}
\isub{a. Réné Descartes}{18}
\isub{b. Thomas Hobbes}{21}
\isub{c. Baruch Spinoza}{26}
\isub{d. Gottfried Wilhelm Leibniz}{32}
\ientry{III. The Philosophy of Experience}{35}
\isub{a. John Locke}{36}
\isub{b. George Berkeley}{38}
\isub{c. David Hume}{39}
\ientry{IV. The Philosophy of the Enlightenment}{43}
\isub{1. Voltaire, Holbach and Rousseau}{44}
\isub{2. Wolff and Lessing}{48}
\ientry{V. Immanuel Kant and the Critical Philosophy}{50}
\isub{Biography and Characterization}{50}
\isub{1. Epistemology}{53}
\isub{2. Ethics}{59}
\isub{3. Theory of Unity on an Aesthetic and Biological Basis}{61}
\ientry{VI. The Philosophy of Romanticism}{63}
\isub{a. Fichte, Schelling and Hegel}{63}
\isub{b. Schopenhauer, Kierkegaard and Feuerbach}{67}
\ientry{VII. Positivism}{74}
\isub{1. Saint Simon and Comte}{75}
\isub{2. Stuart Mill}{77}
\isub{3. Darvin and Spencer}{78}
\ientry{VIII. Recent Treatments of Problems}{81}
\isub{a. The Psychological Problem}{82}
\isub{b. The Problem of Existence}{85}
\isub{c. The Problem of Knowledge}{88}
\isub{d. The Problem of Valuation}{92}
\clearpage

\chapter*{PREFACE TO THE FOURTH EDITION}
\markboth{Preface}{Preface}

In this new edition I have entirely reworked the little book in the light of the
experience I have gained by using it for a number of years in the propaedeutic
course. I have left out much that was seldom needed in that course, but have in
return included biographical information, and I have also tried to make the
presentation more rounded and intelligible. This applies not least to the last
section, which was to give an idea of the present work of philosophy. --- I offer
my best thanks to Dr.\ Anton Thomsen for a number of valuable remarks, which have
benefited this new edition.

\textit{The 28th of March, 1910.}

\hfill\textbf{Harald Høffding.}
\clearpage

\mainmatter
% --- p. 1 ---
\setcounter{footnote}{0}%
\opage{1}\introhead{INTRODUCTION.}

\sectionhead{Philosophical Problems.}

The history of philosophy presents the attempts that thinkers have
made to discuss the ultimate problems of knowledge and of existence.
In modern philosophy, i.e. the philosophy of the last three and a half
centuries, four main problems in particular have been treated,
as the following account will show.

1. \textit{The psychological problem} arises with the question of
the main characteristics of mental life. Is the soul a peculiar substance,
or does its essence consist in a distinctive activity? Is it
composed of many independent elements, or does it
have the character of unity and wholeness? Which of the
various sides or properties of mental life is the fundamental one?
What laws connect the various mental elements
with one another?

Whereas psychological consideration has the factually given
mental life for its object, there are two problems that are conditioned
by the opposition between the factual and the valid
or valuable mental life, --- between the actual constitution of mental life
and the ideal demands that can be made of it:
the problem of knowledge and the problem of valuation. In these two
problems the point is to show and to apply
definite standards of validity or value.

2. \textit{The problem of knowledge} (the logical problem) arises
when the conditions and limits of the validity of our
% --- p. 2 ---
\setcounter{footnote}{0}%
\opage{2}knowledge are asked about. Knowledge is a human mental activity, a work we undertake with our mental powers
and according to the laws of the working of these powers; to what extent are
we then entitled to ascribe validity to its results beyond
our own mind? In what does the truth of knowledge consist? ---
\textit{Formal logic} shows how we must think in order to be
in agreement with ourselves, to be consistent. But the
actual problem of knowledge arises only when we ask
what significance the order and coherence that we establish
among our thoughts have for the subjects that we do not ourselves
produce, --- --- for existence, which is more comprehensive than
our thinking. \textit{Epistemology} (logic in the wider sense)
therefore does not stop at formal consistency or
systematics, but raises the question of the real significance of this
consistency itself.

3. An analogous question is raised by \textit{the problem of valuation} (the
ethical-religious problem) with regard to our expressions of feeling and
will. We judge or evaluate human actions
(our own and those of others) and human social conditions
and call them perfect or imperfect, good or
evil, noble or base. What standard do we have for such
valuation? Is it possible to ground a firm and universally valid
standard and to carry it through with consistency? It is the task of
\textit{ethics} to discuss these questions. But we also evaluate
existence as a whole according to the relation in which it
shows itself to stand to our ethical standard, to our ideals, our
highest values. Does the course of things in nature and in history
show itself favorable or unfavorable, promoting or hindering
our ethical ideals? On the answering of this question
the religious conception (in the widest sense of the word) depends.
It is the task of the \textit{philosophy of religion} to discuss the relation of the
ethical ideals to experience, just as it is the task of epistemology
to investigate the relation of our logical ideals to experience.

4. Finally it can be asked about the nature of the existence
of which the thinking, feeling and willing being
is only a single link.  Thereby \textit{the problem of existence} (the
% --- p. 3 ---
\setcounter{footnote}{0}%
\opage{3}cosmological or metaphysical problem) is posed. What is at issue here is the possibility
of a scientific world-view. Are we able
to comprehend all experiences into a whole, and what validity
can an attempt in this direction have? Every such
attempt will get its distinctive stamp from the domain of experience or the
side of experience that it emphasizes in particular, and in whose light it sees
the other domains or sides. One can thus regard
the mental as the essential, as the primal phenomenon, and
arrive then at an idealist world-view, or make
what is movable in space the primal phenomenon, and arrive then at a
materialist world-view. One can emphasize the unity and
the coherence in the world, and arrive then at a monistic
world-view, or the multiplicity and diversity,
and arrive then at a pluralistic world-view. Within
each of these views the question can again be asked whether existence
is in development, or whether in its innermost ground it is
unchangeable, and whether values are realized in existence,
or whether its essence is exhausted in a system of causes and
effects. This last question leads to a connection of the problem of existence
and the problem of valuation.

Whereas history shows that there has been a tendency,
by reason of a sanguine confidence in the power of thought, to go immediately
to the problem of existence, there appears later a decided
endeavor to treat psychology, epistemology and
ethics independently of ``metaphysics'' or the systematic treatment
of the problem of existence. But the interest in philosophy
can spring from four different motives, corresponding to the four
main problems, and the one problem leads naturally over to
the other, in so far as there is an analogy between them. One and the
same fundamental problem runs through all of them: the relation
between unity and multiplicity, continuity and discontinuity
comes forward under different forms, whether we investigate
personality, science, the conduct of life or the world-view.
The spiritual struggles in these different domains
present analogous features: what is at issue is the possibility of
comprehending into a unity what provisionally appears as
a multiplicity, which offers greater or lesser resistance to
% --- p. 4 ---
\setcounter{footnote}{0}%
\opage{4}giving up its sporadic character. All the main problems consist
in tasks for psychical energy. If a conclusive
solution of them is unattainable, this means only that however far
we may get in our inquiry, new horizons constantly form
for us, with new goals and new tasks\footnote[1]{Cf. the essay \textit{Filosofien og Livet} in my ``Mindre Arbejder''. Second Series.}.

% --- p. 5 ---
\setcounter{footnote}{0}%

\closerule

\opage{5}\chapterhead{I.}{THE ORIGIN OF MODERN PHILOSOPHY.}

\sectionhead{1. The Discovery of Natural Man.}

a. \textsc{Burckardt} in his work \textit{Die Kultur der Renaissance
in Italien} has characterized the Italian Renaissance as a
discovery of man. The historical conditions led to a liberation
of individuals. Man was no longer judged merely
by his relation to the Church or to his estate (corporation),
but became an independent object of interest and investigation.
The discovery of ancient art and literature contributed to this.
Men learned that outside the Church too there was to be found a distinctive
human life with laws and purposes of its own.
The horizon was widened, and a comparative consideration
became possible.

In Italy itself \textsc{Niccolo da Machiavelli} (1469---1527)
can be regarded as a characteristic representative of these endeavors.
He belonged to an old Florentine family and entered
the service of his native city as a diplomat, whereby he had
good opportunity to observe men and events. After
the fall of the Republic (1512) he entered into relations with the Medici,
but thereby aroused the ill will of his fellow citizens, and when
the republican constitution was introduced again, his services were not
accepted. He died in dejection in 1527. --- It was
political interest that made him a thinker. The unhappy
conditions of Italy awakened in him the wish to see the strength and way of thinking of
antiquity awakened again. Why should one only
imitate the art of antiquity and not its deeds? But
% --- p. 6 ---
\setcounter{footnote}{0}%
\opage{6}great deeds are possible only when one ruthlessly seeks to
reach the goal one has set oneself! --- --- It may sometimes seem in Machiavelli
as if it were a matter of indifference what goal
one sets oneself, if only one pursues it with ruthless
energy. But there was one goal that always stood in the background
of his thinking: the unity and greatness of Italy. In the service of this
goal all means were sacred to him. In his famous
book Il \textit{Principe} (1513) he has expressed himself in this direction.
When he found no sense for the great and no
true energy in the Italians of his time, he derived it from the fact that the Church
and Christianity had weakened the passions. In his
\textit{Discorsi} (Treatises on the First Ten Books of Livy) he compares
the way of thinking of antiquity and of modern times, and
the comparison does not turn out to the advantage of the latter.
Honor, strength of mind and bodily strength are not sufficiently valued.
In his admiration for these qualities Machiavelli
is a true son of the Renaissance.

In France the spirit of the Renaissance appears in \textsc{Michel} \textsc{de}
\textsc{Montaigne} (1533---1592), a nobleman who lived a rich
literary life at his castle near Bordeaux, far from
the political and religious struggles of the time. His interest was the study
of human life, as he knew it from travel, books
and self-observation. His \textit{Essais} (1580---1588) begin with
the statement: Je suis moy mesme le sujet de mon livre.
Our great mother Nature reveals herself in each individual in a distinctive
way. The more we investigate, the more differences
we discover, without being able to bring them under general
laws. Therefore Montaigne speaks against all superstitious faith in
reason (rationalism) as well as against belief in dogma. We know
things only as our senses show them to us, and in all
notions of God we clothe him with human attributes,
just as animals, if they could think, would picture their
gods in the shape of animals.

In purely theoretical form the interest in the human
appears in the learned Spaniard Luis Vives (b. in Valencia
1492, d. in Bruges 1540), who by his book \textit{De anima et vita}
(1538) is a forerunner of modern empirical psychology. He
% --- p. 7 ---
\setcounter{footnote}{0}%
\opage{7}demands that experience shall be the foundation of all knowledge,
and by virtue of this principle he declares that our main interest
is not to know what the soul is, but \textit{how} it works.
His method is descriptive, not analytic and explanatory;
but his description of the individual mental phenomena,
especially of the emotions, is still of interest. For him the soul and
the vital force are One, and he seeks everywhere to connect
physiology and psychology.

b. Within Protestantism \textsc{Melanchton} (1495
--- 1560) worked for philosophy through a series of textbooks. But
the greatest significance was attained by Calvinism, within which the struggle
for religious freedom brought about endeavors toward a liberal
doctrine of law and of the state. Thus the burgomaster of Emden
\textsc{Johannes} \textsc{Althaus} (Althusius) (1557---1638) developed the
idea of the sovereignty of the people in his book \textit{Politica methodice digesta} (1603). But
most famous in this domain is the Dutchman \textsc{Hugo} \textsc{de}
\textsc{Groot} (Grotius) (1583---1645). Early he attracted attention
by his juridical and theological learning. As a politician
he belonged to Oldenbarnewelt's aristocratic and liberal-theological
party. After Oldenbarnewelt's death he was imprisoned,
but escaped by the ruse of his wife and afterwards lived in Paris,
finally as Swedish ambassador. In his work \textit{De jure belli et pacis}
(1625) he proceeds from the fact that there is war in the world --- 1) between
states, 2) between an individual and the state, 3) between individuals
among themselves, and 4) between the state and an individual. 1) The rights of man
and the duties of humanity are not abolished in war, which
is to be waged in a human manner and for the sake of peace. By
this principle Grotius has become the founder of modern
international law. 2) The individual's struggle against the state is rebellion, which
Grotius rejects (whereas Althaus maintained that it could
be justified). 3) War between individuals ought to be limited to
self-defense. 4) The state's war against the individual appears as
punishment. This is justified only when the suffering that is inflicted
on the lawbreaker is outweighed by a greater good for himself
and for society. --- In all these relations right holds independently
of theological grounds, springs from principles that
lie in the nature of man (ex principiis homini internis),
% --- p. 8 ---
\setcounter{footnote}{0}%
\opage{8}namely as the condition of the social life for which man
has an involuntary need. The most important condition here is that
every promise be kept; to this men have bound themselves
by an express or tacit contract, as soon as they agree
to live in a society.

At about the same time that Grotius gave natural law and international law
a purely human foundation, the English
diplomat \textsc{Herbert} \textsc{of} \textsc{Cherbury} (1582---1648) maintained a natural religion
springing from the nature of man himself \textit{(De veritate}
1624). He opposes both those who place an ecclesiastical faith
above knowledge of reason, and those who trust only
the investigating reason, and those who would derive everything from
the senses. There is, he maintains, an immediate, instinctive
sense which leads men to the recognition of certain truths
(notitiæ communes), a sense that is an effect of the
instinct of self-preservation implanted by God. Such truths
are e.g. that two propositions which conflict with each other
cannot both be true, --- that there is a first cause
of everything, --- that we should do to others what one wishes that they
should do to us. Natural religion also rests on
an instinctive foundation, an inner revelation in every human soul.
A testimony of this he finds in the fact that we have faculties
and feel a need that cannot be satisfied by any
finite object. The content of natural religion is the belief
in a supreme being, who is to be worshipped by virtue and pious
disposition, and in reward and punishment after death. Questions that
go beyond this need not be answered.

In the domain of the theory of knowledge another English
nobleman sought at the same time to lay a new foundation. \textsc{Francis} \textsc{Bacon}
of Verulam (1561---1626) expected a mighty rise in culture
when men, by the aid of natural science, which
antiquity and the Middle Ages had neglected, made themselves masters
of the forces of nature. The purpose of science is the enrichment of
human life by new inventions. Yet Bacon also emphasizes
the joy itself in the contemplation of nature (contemplatio
rerum): the sight of light is more glorious than the manifold
utility of light! --- In order to carry out his great scientific
% --- p. 9 ---
\setcounter{footnote}{0}%
\opage{9}plans he considered himself entitled to use all means.
What mattered was power and wealth for the costly investigations.
Therefore he plunged into political life and
reached --- by compliance with the despotic whims of Elizabeth and James the First --- the highest
positions. As Lord Chancellor
he was accused by Parliament of bribery and
other illegality and removed from office. He spent his last years in
quiet with his works.

Bacon calls himself a herald who proclaims a new
era but does not himself fight along. He demands that one should
give up fruitless speculations and apply the method of experience
(induction) in all domains, the mental as well as the
material. Against the induction that had earlier been applied
Bacon objects that one has kept only to the positive
cases, and he demands that one should also investigate the cases
in which a phenomenon does not occur. Further he demands
an investigation of the manner in which a phenomenon
varies under different conditions, and he lays stress on
the point that one should --- in order not to be overwhelmed by one's material --- try
to set up hypotheses and see whether they are confirmed by new experiences.
Yet he does not --- because he underrates the quantitative
method --- arrive at the true method of the new natural science,
which we shall find in his contemporaries Kepler and
Galilei. --- It is in his work \textit{Novum Organum} (1620) that he
describes the inductive method, after first having investigated
the grounds of the imperfection of the sciences.

In order rightly to grasp nature man must make himself a
blank slate; only as a child can one enter into
the kingdom of nature. But there are natural and acquired illusions
(idola mentis) which hinder us in this. They can be traced to
four kinds. The first kind is common to the whole species, because
it lies in human nature (idola tribus). We are inclined to grasp things in likeness to ourselves (ex
analogia hominis) instead of according to their place in the connection of the world
(ex analogia universi), and we are inclined to
find order, simplicity and regularity in nature. The
second kind of illusions rests on the individuality of the individual
% --- p. 10 ---
\setcounter{footnote}{0}%
\opage{10}(idola specus: Everyone looks out on nature from his own cave).
Some emphasize most the similarities, others most the differences;
some emphasize the wholes, others the elements of which they
consist. The third kind is due to the influence of language on thinking
(idola fori). Words are formed according to the needs of
practical life and often indicate differences or combinations that science
cannot approve. The fourth kind (idola theatri)
is due to the influence of traditional theories. --- Bacon does not investigate
how it is possible to cast off these illusions.
In particular there lies in idola tribus a great problem, which in a later
time was brought forward: we must always regard existence from
a human point of view (ex analogia hominis), but how
can our knowledge then have absolute validity?

\sectionhead{2. The New Picture of the World.}

In its view of nature, as in its conception of man,
the Middle Ages, in so far as it did not take its ideas from
the Bible and the Christian tradition, had built on Greek antiquity.
Its medicine it learned from Galen, its astronomy from Ptolemy,
its philosophy from Aristotle. Its world-picture was due to
the combination of the Ptolemaic and Aristotelian doctrine with
the biblical conception. The earth stood fixed in the center of the world.
Around it revolve, fastened on solid but transparent
spherical surfaces (spheres), the moon, the sun, the planets and the fixed stars
(these last all on one and the same sphere). The earth and
the space from the earth to the moon, the sublunary world, is
the home of changeability and perishability; here the
four elements (earth, water, air and fire) are in constant motion,
as they seek their ``natural places,'' the heavy ``downward,'' the
light ``upward.'' Above the sphere of the moon is found the ether (the fifth
element), a matter that has no ``natural place'' and therefore
continues its motion eternally and with absolute regularity.
By this unchangeability the motions of the heavenly bodies are
an image of the essence of the Deity. They move
in circles; the circle is the most perfect figure, because
it always returns into itself. The world is bounded by
% --- p. 11 ---
\setcounter{footnote}{0}%
\opage{11}the sphere of the fixed stars, which is moved immediately by the Deity,
whereas the lower spheres are moved by subordinate star-spirits.
--- This picture of the world agreed not only with what
the authorities of the time, the Bible and Aristotle, taught, but also
with what the senses immediately seemed to teach. It cost
therefore a hard struggle to displace it. And this struggle
brought with it a remodeling of the whole world-view.

The founder of the new world-picture, Nicolaus Copernicus
(Coppernick), was born in Thorn (1473), studied in
Cracow and at various Italian universities and lived
as a canon in Frauenburg, occupied partly with studies,
partly with administration. In the great struggle of the age he took no
part, although he seems to have had a certain sympathy with the
religious movement as well as with humanism. As early as
1506 he began the elaboration of his new astronomical
theory, but hesitated long to publish his work \textit{(De
revolutionibus orbium coelestium),} and only on his deathbed
(1543) did he get a printed copy into his hands. --- Here we are interested
especially in the epistemological presuppositions on which
his work rests.

Nature always works in the simplest way. Against this
it would conflict if the whole world moved around so small a
body as the earth. And in addition, in order to explain
that the planets did not, like the fixed stars, always move in
the same direction, but apparently back and forth, one had
been obliged to assume that they did not move in simple circles, but
in small circles which (as epicycles) moved on the
larger circles. A far simpler picture of the world is obtained when
the sun is thought of as the center of the world, and the earth and the planets
move around it.

Besides the principle of simplicity, the thought of the relativity of
motion plays an essential role in Copernicus. An observed
motion can be explained not only by the observed
object's moving, but also by the observing
subject's itself moving. If it is now assumed that the earth, from
which the motions of the heavenly bodies are observed, is itself in
motion (about its axis and around the sun), the observations could
% --- p. 12 ---
\setcounter{footnote}{0}%
\opage{12}be explained just as well, but far more simply, than according to
the old system.

Copernicus still assumed the world to be bounded by the sphere of the fixed stars,
which he thought was the immovable boundary of the world, and
he believed in solid spheres in which the stars were set.
But he had nevertheless made a decisive change and
had stressed that we cannot infer with certainty from our conception
of existence to the nature of existence itself.

For long Copernicus' doctrine was regarded as an idle
thought experiment, and almost half a century passed
before it won adherence among independent investigators.

\textsc{Giordano} \textsc{Bruno} stands here as the one who not only drew
further-reaching consequences of the new doctrine, but laid
it as the foundation of a whole world-view, in its main features
the distinctive fundamental conception of modern times and its
first attempt at an independent philosophy. Bruno was born in
Nola in southern Italy (1548) and entered the
Dominican order while quite young. When he made himself guilty of heresies,
and when his need to move freely made the constraint of the cloister
unendurable, he fled, cast off the monk's habit and began
a restless life of study and wandering, which led him to Switzerland,
France, England and Germany. He appeared as a teacher in
Toulouse, Paris, Oxford and Wittenberg, but found no
lasting place, partly because of the resistance of the old schools,
partly because of his own restless and passionate spirit.
Yet he found time to compose his ingenious works, of which
especially the Italian dialogues published in London in 1584 (which
are now available in German translations) are to be singled out. When he
considered a reconciliation with the Catholic Church possible, and when
he had greater ill will toward the Protestant Church than
toward the Catholic, he returned to Italy, where he wished to
live in peace for his studies. But he was imprisoned by the Inquisition
and after a long captivity burned in Rome in the year
1600.

Bruno praises Copernicus for the lofty and free spirit that
had raised him above the testimony of the senses, by which the
great multitude remains held. But he disputes that the world
% --- p. 13 ---
\setcounter{footnote}{0}%
\opage{13}can end with the sphere of the fixed stars, and supports himself herein partly on
epistemological, partly on grounds of the philosophy of religion.

Only apparently do the senses bear witness to an absolute center of the world
and an absolute whole of the world. Every time we change standpoint,
we get a new center of the world and a new boundary of the world,
so that each place in the world can be regarded now as center,
now as boundary point. And it goes with thinking
as with sense perception: number can be added to number and thought to thought
without limit. Every place is determined by its relation to
other places, as every number by its relation to other numbers and
every thought by its relation to other thoughts. Every motion
looks different according to the standpoint in relation
to which it is conceived. Since time is measured by motion,
the determination of time too becomes always relative. Neither
are heaviness and lightness now absolute properties, but
hold in relation to the larger bodies to which the elements belong.
The old opposition between heaven and earth must also
fall away. We have no right to assume that the fixed stars are each
by itself unmoved and all at the same distance from us; it
might well go with us here as when distant ships seem to us to be
immovable and all equally far away. It is in any case
something that must be investigated more closely. And the fixed spheres must be
rejected, because Tycho Brahe has shown that the comets are not
(as was formerly believed) meteors in the earth's atmosphere, but
pass straight through the heavenly space and thus long ago would have had to
break the spheres. The stars must then also, just as well as
the earth, be able to move without support.

From the standpoint of the philosophy of religion Bruno proceeds from the infinity of the Deity.
The effect must be infinite if the cause is.
The divine fullness could not come to unfolding
in a bounded world.

As it goes with the opposition between heaven and earth,
so, according to Bruno, it goes with all oppositions with which we
work, because we can think only by setting the one over against
the other. Neither the opposition between matter and form,
nor that between possibility and actuality, nor that between
matter and mind may we ascribe absolute validity. In the essence
% --- p. 14 ---
\setcounter{footnote}{0}%
\opage{14}that lies at the ground of everything, but for which no positive
determination of thought can be found, all these oppositions must
fall away. The Deity cannot have its seat in
``heaven,'' since the opposition between heaven and earth does not
hold, but every point can be seen both as center and
as boundary point. It works in our innermost being, in our striving
for self-assertion; it is the soul in our soul, as it is
the soul in all of nature, and through space it brings everything
existing into interaction. But it is honored better by silence
than by speech.

Bruno's ethics stands in close connection with his world-view.
Among the virtues he puts love of truth and
zeal for work in the first rank. No goal can be the last,
but every goal begets new striving. Where a
heroic enthusiasm (furore eroico) has been aroused, man does not let himself be
deterred from striving after a high goal because pain
and danger are bound up with it; they are evils only for the sensuous
consideration. The higher the goal, the greater the pain
that is possible. But the heroic man rejoices when the noble
fire is kindled in him, even if the goal should not be reached.

\sectionhead{3. The New Science.}

The struggle over the validity of the new picture of the world could be
brought to an end only when it could be demonstrated as a necessary
consequence of physical laws of nature. This happened through \textsc{Isaac} \textsc{Newton}'s
demonstration of gravity as a universal law (1687). But the presupposition
for this was in turn a new scientific method; it is
to such a method that modern natural science owes its origin.
Instead of speculation and construction on the one hand and
mere gathering of observations on the other, there now came
exact measurement, quantitative determination of the phenomena, and,
on the basis of it, strict calculation and inference, whose results
could in turn be confirmed by new measurable phenomena.
Modern natural science came into being through such an interaction
between induction and deduction, observation and mathematical
calculation. It is Kepler and Galilei to whom the introduction of
% --- p. 15 ---
\setcounter{footnote}{0}%
\opage{15}this method is due, after the famous artist
\textsc{Leonardo} \textsc{da} \textsc{Vinci} (1452---1519), the most many-sided
genius of the Renaissance, had already expressed thoughts in the same direction, which however became known only
much later.

In \textsc{Johannes} \textsc{Kepler} (1571---1630), the famous astronomer,
there took place an interesting development from a theological-poetic
world-view to an exact natural-scientific conception.
He first proceeded from the view that the universe, as a revelation of
the Deity, must be ordered according to simple geometrical
relations, which he even thought he could construct in advance.
Later he held fast only to the general conviction
that certain mathematical laws must hold for the motions
of the planets, and on the basis of Tycho Brahe's
observations he found these laws (having shown, among others, that
the orbit of Mars was an ellipse). And whereas earlier he had
explained the motion of the planets by the assumption of peculiar
planetary souls, he later (in his writing on the orbit of Mars) set up
the principle that the causes one assumed should themselves be capable of
being shown in experience; such a cause he called \textit{vera
causa.} The planetary souls were not veræ causæ. --- Besides the
theological grounding of the quantitative method he has
also a psychological grounding: mathematical knowledge
is the clearest and surest that we possess, and ought therefore to be
applied as far as possible. In qualitative respect the
phenomena may present themselves differently to different individuals;
but with quantitative determination the disagreement ceases. ---
Finally experience shows that all material phenomena have
quantitative, especially geometrical properties: ubi materia, ibi
geometria.

The real founder of modern natural science
is, however, \textsc{Galileo} \textsc{Galilei} (1564---1642), through his clear
consciousness of the interaction between induction and deduction,
observation and calculation. Observation cannot embrace
all possible cases; we must therefore keep to certain selected
phenomena, describe and measure them as exactly as possible,
on the basis of this set up by inference and calculation a
hypothesis, and investigate whether the consequences that can be
% --- p. 16 ---
\setcounter{footnote}{0}%
\opage{16}derived from it agree with what observation shows. Both
in the setting up and in the testing of the hypothesis quantitative
determinations are necessary, and Galilei therefore set up the
principle: measure everything that is measurable, and make measurable what
is not yet so!

Kepler had already taught that a body cannot of itself
pass over from rest to motion. Galilei goes a step
further. It seemed to him the simplest (and like Copernicus,
Bruno and Kepler he regarded the law of simplicity as a
universal law) that a thing remains in the state in which it
is, when no external causes intervene. A body can therefore
not of itself change its motion. If all external
influences could be kept away, the motion would continue constantly
with the same velocity. By experiments Galilei demonstrated
thereupon that balls rolled along parchment channels continued
their course the longer, the smoother the channels were, and the smoother
and more regular they themselves were. Circular motion too,
which he regarded as simple, would in his opinion continue constantly
in the absence of external hindrances, --- not merely
rectilinear motion.

These physical theories Galilei developed in \textit{Discorsi delle
nuove scienze} (1638). Some years earlier he had in his
\textit{Dialogo sopra i due massimi sistemi del mondo} (1632) compared
the Ptolemaic and the Copernican system of the world
with one another. He thought he had weighed the possibilities
for and against without pronouncing decisively; but his own
view was clear, not least from the irony with which
the representative of the old view was treated. This
brought about the catastrophe of his life. Already earlier (after the discovery
of the moons of Jupiter and of the sunspots) he had publicly
declared himself in favor of the Copernican system. And
when in the year 1616 the Inquisition put Copernicus' book on the list
of forbidden books, Galilei is said to have promised not to hold
or to spread the new doctrine. The Inquisition saw in his Dialogue
a breach of this promise, and the seventy-year-old man
had to, under threat of the rack, abjure the impious
doctrine that the earth moved and was not the center of the world.
% --- p. 17 ---
\setcounter{footnote}{0}%
\opage{17}Until his death he stood under the surveillance of the Inquisition, and his
later works had to be printed secretly abroad.

Galilei stresses strongly that Copernicanism enjoins
that we must not confuse our conception with existence itself:
``Think the earth away, and the sun does not rise and does not go
down, and there is neither horizon nor meridian, neither
day nor night!'' This thought he extends to the whole sensuous
view of nature under the influence of the conviction
that we obtain true knowledge only by keeping to the
quantitative properties of things. Only size, figure and motion
or rest can we ascribe to things themselves; they are the first
and real properties. On the other hand it is a sensuous prejudice to
regard taste, smell, color, heat etc. as absolute properties
of things. They are names that we give to things
when they call forth certain sensations in us; they have their
seat in our sensuous body and fall away when this falls
away. This doctrine, in which the principle of mechanical natural science
is expressed, became of great importance for the epistemological
investigations of the following age.

\closerule

\chapterhead{II.}{THE FORMATION OF SYSTEMS.}

The new interests, points of view, and discoveries that the age of the
Renaissance had brought were bound to call forth an effort
to develop, on the foundation thus given, a
new world-view coherent in itself. There was bound
to stir a need to draw the ultimate consequences of the
new views. At all times there stirs, more or less
distinctly, a need in the human mind to bring together
its knowledge in a view of the whole. The task now presented itself
in particular of uniting the new world-picture and the new
% --- p. 18 ---
\setcounter{footnote}{0}%
\opage{18}science with the remaining content of the life of the spirit. Bruno had already
worked in this direction; but the new mechanical view of nature,
through which many problems took on a different shape
from before, could not influence his way of posing problems.

a. \textsc{Réné} \textsc{Descartes} (1596---1650) may be called the true
founder of modern philosophy, because he is the first to seek out methodically
the presuppositions of knowledge, and because he was the
first in whom the mechanical explanation of nature acquired a decisive
influence on the way problems were posed. He was the son of a
French nobleman, and his economic independence made it
possible for him to devote his life to thinking and inquiry. In
his \textit{Discours de la méthode} (1637) he has given an interesting
philosophical autobiography. He received his first instruction
at a Jesuit college. Although he had the best
teachers of the age, he was nevertheless not satisfied with the result
of his schooling. He knew a great deal, but coherence
and clear grounding were lacking. He then threw himself
into the life of the world and resolved later to get to know real
life through military service and travel. Yet he was
again and again drawn back to philosophy. When in the year 1619
he lay in winter quarters in Bavaria, he had a kind of scientific
awakening. The path he was to take as an investigator stood suddenly
before him: The bewildering multiplicity of our representations
will vanish when we begin with the simplest and clearest
thoughts and then proceed step by step to the composite
questions, constantly ordering our thoughts in such a way that
the later can be derived from the earlier. This idea guided
him both in his mathematical works (the founding
of analytic geometry) and in the working out of his
philosophical system. After some years of study in Paris he went to
Holland in order to live in greater peace, and because in France
strict prohibitions were issued against departure from the medieval
philosophy. But as the Dutch theologians no less than the
Catholic ones regarded his philosophy as dangerous, he accepted an
% --- p. 19 ---
\setcounter{footnote}{0}%
\opage{19}invitation from Queen Christina and spent his last years in
Stockholm.

In order to find the first principles of knowledge,
Descartes (most fully in the \textit{Meditationes} 1641) examines whether there is
anything that cannot be called in doubt. Of every object,
every content of our thoughts we can doubt. Every observation
and every assumption can be illusion. But even the
most radical doubt is a thinking, an act of consciousness; if
all thinking falls away, then all doubt falls away too. If I doubt,
then I still know that I doubt. Thought itself I can therefore
not doubt, even though I doubt all its results.
The word thinking Descartes takes in a broad sense, of every kind of
consciousness, feeling and willing. That I think shows that I
am (as thinking): je pense, donc je suis (cogito, ergo
sum)! This is, however, no inference, but an immediate consciousness
(of being conscious). --- The clearness and distinctness
with which I know that I think is now for Descartes a
standard: all propositions that are as evident as
this one must be valid. And there are two important propositions that
satisfy this demand: the proposition that everything has a cause,
and the one that there cannot be more contained in the effect than in
the cause. From these propositions Descartes derives a proof of
God's existence: whereas all representations of finite things
and beings can find their explanation in us or outside
us, this is not possible for the representation of an infinite
being; it can be due only to an infinite cause. Moreover
it would be a self-contradiction if God's existence were rejected:
for the concept of God is the concept of a perfect
being, and to perfection existence also belongs.
It follows, then, from the concept of God that he must exist.

Only now, when Descartes (admittedly with the help of the causal
principle) believes he has proved the existence of God, is the truth of our knowledge
secured for him; for a perfect being cannot
deceive. Like Kepler, then, he gives his doctrine a
theological foundation. Yet the concept of God with which he operates
here is not a purely religious concept. It expresses
the principle of the lawful connection in nature. Everything
% --- p. 20 ---
\setcounter{footnote}{0}%
\opage{20}that is to be recognized as true must find its place in
this connection, which makes it possible for us to distinguish between
dream and reality.

By means of the concept of God Descartes then grounds
the first principles of natural science. \textit{(Principia} Philosophiæ.
1644). Only three properties have we a right to ascribe to
material things: extension, divisibility, and mobility.
And we obtain the clearest knowledge of everything that takes place
in the material world when we confine ourselves to these three properties;
the sense qualities, on the other hand, must be charged to the account of the perceiving
subject. Herein Descartes pronounces the
fundamental thought of mechanical natural science, which he seems to
have arrived at independently of Galileo, but probably under
the influence of Kepler. --- From God's immutability (which is
a part of his perfection) it follows now that a body remains
in its state when an external cause does not act upon it.
And not only the principle of inertia is thus grounded theologically.
From the divine immutability there follows further the principle
of the conservation of motion (an imperfect precursor
of the principle of the conservation of energy): the sum of motion placed
in matter at creation can be distributed differently
at different times, but no motion is lost, and no
wholly new motion arises.

Proceeding from these presuppositions Descartes demands that everything in nature be
explained in a mechanical way. He himself has satisfied
this demand only through hypotheses that have not always
stood the test. --- --- Like the world as a whole, so too the
organism is a machine. A confirmation of this view
Descartes found in Harvey's discovery of the circulation
of the blood, which showed that the motion of the blood was due to mechanical
causes. He himself applied a mechanical explanation to the
nervous system, in that he describes the involuntary movement
now called reflex movement, in which an impression,
without any consciousness needing to be present, immediately
releases muscular activity. Consciousness he assumed only
in man, whereas he regarded animals as mere
machines.

% --- p. 21 ---
\setcounter{footnote}{0}%
\opage{21}The soul of man stands in reciprocal action with the brain, in particular
with a definite part of it (the pineal gland, glandula pinealis), which
in his opinion lay centrally in the brain and which was not, like
the other parts of the brain, double. Upon this gland the ``animal spirits''
(the fine air that, according to the physiology of the age, derived from
antiquity, flowed through the nerves) impinge, and the impact
is then propagated to the soul, whereby sense-sensations
arise. When the soul in its turn impinges on the ``gland,'' a
change is produced in the direction of the ``animal spirits,'' which in turn causes a
muscular movement, that is, the execution of an act of will. By this clear and consistent carrying through of
the popular notion of a reciprocal action between soul
and body Descartes comes into conflict with his own doctrine of the conservation of motion:
motion disappears when the pineal gland impinges on
the soul, and motion arises when the soul impinges on
the pineal gland. Even though the soul only changes the direction of motion,
this conflicts with the principle of inertia.

Soul and body are for Descartes two different beings
(substances), because they present different properties, thinking
(consciousness) and extension, which cannot be derived
from one another. How two such different beings could
act upon one another became accordingly a difficulty for Descartes
and his followers. ---

Cartesianism was the first form in which the thought of
the new age became accessible to wider circles. By its
strict demand for a mechanical explanation of nature it marked,
despite its often unfortunate hypotheses, a great advance, and
it aroused a lively movement in the domain of natural science.
By its spiritualism (the assumption of the soul as a being different from the
body) and by its attempt to combine
theology and science it could --- despite violent opposition ---
find entrance into religious circles. By its great clarity
his doctrine could easily be presented in a popular way, and it won wide support
in France, Holland, and Germany.

b. \textsc{Thomas Hobbes} (1588---1679) made the first independent
attempt to regard the mechanical natural philosophy
% --- p. 22 ---
\setcounter{footnote}{0}%
\opage{22}as the only science, --- its points of view as the
only ones from which existence can and should be conceived. This man,
energetic as a thinker and polemicist, but mild and timid in his life,
was, like Descartes, dissatisfied with the medieval
school he had passed through. For a time he occupied
himself mostly with literary studies. The unrest in
England awakened his interest in political and ethical questions,
and when he had come to know the new scientific
points of view (especially after he had become, on a journey in Italy,
a personal friend of Galileo) he was led over to general
philosophical studies. He lived as a private man with
the noble family Cavendish, first as a tutor, afterwards
as a friend of the house.

When the Revolution broke out, Hobbes went to France, where
he was for a time later the tutor of the exiled Charles II.
When Cromwell had brought peace to the country and granted freedom of expression,
he returned. ---

Hobbes seems to have arrived by his own path at the
conviction that all change is motion. Being
present at a discussion about what sensation is, he hit
upon the thought that if everything in nature were immovable
or in uniform motion, no sensation would arise.
A difference of motion (diversitas motuum) is therefore a
condition of sensation. To sense one and the
same thing constantly and not to sense at all come to the same thing.
On this proposition Hobbes grounded his psychology. When he
then also thought about what change was, he found that
it consisted only in motion, and from this it followed in turn that the sense qualities
(color, sound, etc.) have a purely subjective validity.
He also appeals here to experience, e.g.
to the fact that a blow on the eye produces light, although there is in fact
no external light.

The proposition mentioned, that all change consists in motion,
is for Hobbes the main principle of science. As such
it cannot be proved, but it will be acknowledged when one merely lays aside
all prejudices. In a similar way he sets up a
series of other fundamental propositions: the causal principle,
% --- p. 23 ---
\setcounter{footnote}{0}%
\opage{23}the principle of inertia, the proposition that motion can arise only from motion
and produce nothing but motion. When now these
fundamental propositions are to explain everything, then everything must consist in motion.
In his systematic exposition Hobbes begins
with the doctrine of body \textit{(corpus)} in general, an exposition
of the geometrical, mechanical, and physical laws of motion,
passes then to the doctrine of man \textit{(homo),}
i.e. of the motions that take place in man, the physiological
and psychological motions, and ends with the doctrine of
the citizen \textit{(civis),} i.e. of the motions in men that
condition their mutual relations and their social life. --- The
freshest presentation of his psychological and social theories
he has given in a work written about 1640, \textit{Elements of
Law} (new edition by Tønnies 1888). Later he published partly
the political-philosophical works \textit{De cive} and \textit{Leviathan,} partly a systematic
exposition of the first two parts of the system (De corpore
and De homine) under the name of \textit{Elementa Philosophiæ} (1655).

The first fundamental propositions must, according to Hobbes, be set up
as arbitrary definitions. But he is clear that science
cannot content itself with deductions from such
definitions. When we pass from geometry to mechanics, from there
to physics and from there in turn to psychology and the theory of society, we take up
constantly new presuppositions and set up hypotheses that make possible
the application of the general fundamental propositions to the special
domains. But he does not doubt that the fundamental propositions
can be applied throughout. When everything is motion,
everything must be corporeal. An incorporeal thing is a non-thing. Science
can investigate only finite things, since only such
can be in motion from one place to another. Of the world
as a completed whole we can have no knowledge. All questions
concerning the world as a whole lead over to the inconceivable
and can be decided only by faith, not by knowledge.
Neither about the origin, limit, nor end of the world can science
teach anything. The highest science, the first fruits of wisdom
(primitiæ sapientiæ), is reserved to the theologians, just as
in Israel the first fruits of the harvest were reserved to the priests. ---

When everything is motion, sensation, and
% --- p. 24 ---
\setcounter{footnote}{0}%
\opage{24}consciousness
in general, must also be motion: ``Sensation is
\textit{nothing else than} a motion in the minute parts of the sensing body'';
``Consciousness (mens) is \textit{nothing else than} a motion
in some parts of an organic body.'' Motion is
the real; consciousness is only the appearance (apparition) of motion
to us. Thus the feeling of pleasure is in reality
(really) only a motion in the heart, thinking only a motion
in the head. Psychology is therefore a part of the general
doctrine of motion. Yet Hobbes distinguishes expressly
motion itself from its appearance or apparition
to us. He even declares that of all phenomena the most remarkable
is this, that anything at all can become a phenomenon, that is,
appear to a subject. There must therefore after all be something other than
motion in the world.

In his special psychology Hobbes emphasizes the great
importance that comparison has, already as regards sense perception
itself, and he shows how representations arise
and are connected according to laws of association. He has laid the foundation of
the English psychological school. --- Of particular interest is
his doctrine of feelings and passions. Pleasure or pain arise
when a change furthers or hinders our striving
for self-preservation. An absolute goal that could be attained
once and for all is unthinkable; all striving and feeling would
then fall away. The highest good must consist in unhindered progress
toward ever higher purposes. For life is motion! ---
The special forms of striving and feeling develop from
the feeling of power and of powerlessness. The notion of the future
plays, namely, a great role: if I believe I have means for future
satisfaction and for overcoming obstacles, the feeling
of power arises, in the opposite case the feeling of powerlessness. Among the means
on which I reckon may also be help from friends or from
God. Since my striving for self-preservation is easily crossed by
the striving of others, comparison with them plays a great
role. Life is a race. We exult when we get ahead of
others, feel humiliation when we fall behind, --- --- hope
when we are in good spirits, despair when we are tired, anger
when we meet a sudden obstacle, pride when
% --- p. 25 ---
\setcounter{footnote}{0}%
\opage{25}obstacles are overcome; --- we laugh when we see others fall, weep when
we ourselves fall; --- we feel pity when we see one whom we wish
well fall behind, love when we help others. Only with
death does the race end.

In this great race it is learned that the efforts at self-preservation
of different men are not in mutual
harmony. Strife arises, and encroachment may constantly be feared.
The state of nature, i.e. the condition of human life as it
would be without a political order, is a war of all against all
(bellum omnium contra omnes). Here only the need
and power of individuals count, as far as they can extend themselves. Fear, hatred, and
restless lust for power rule. But in calm moments (sedato
animo) men will realize that they can gain greater
advantages by union and community than by strife.
Thereby they discover the fundamental moral law: Peace is to be
upheld, and if it cannot be attained, war is to be waged in common!
From this fundamental proposition the particular virtues and
duties can be derived, for faithfulness to one's word, gratitude, helpfulness,
forbearance, fairness, and self-control are necessary
qualities if peace is to rule and society to
subsist. Likewise there follows from that fundamental proposition the rule that
one must not do to others what one does not wish that they
should do to us.

But for the practical carrying out and upholding of these laws
and rules a strong power of the state is necessary; the freedom
of the state of nature must be given up. This takes place by an express or tacit
contract, by which all individuals at one stroke
renounce their freedom and bind themselves to unconditional
obedience to an authority (a prince or an assembly).
Only by such an absolute submission can peace be secured.
From now on the will of the government becomes one with the will of the people.
Hobbes is the naturalistic theorist of absolute sovereignty.
Every limitation (by estates or Church) and every
division (between prince and parliament) would contain the seed
of war and could lead back to the state of nature. In the
moral and religious domain too the will of the sovereign is to
be decisive. Above all the sovereign is to decide how
% --- p. 26 ---
\setcounter{footnote}{0}%
\opage{26}God is to be worshipped; otherwise what is worship for one
would be blasphemy for another, and there would
be a constant source of strife and dissolution. For the same reason
the first definitions of good and evil must be fixed by the sovereign.
Hobbes here constantly presupposes that the sovereign
expresses the will of the people to hold the state of nature down and
uphold peace. He is a precursor of the enlightened
despotism of the eighteenth century. What he fought was first
and foremost class rule and the domination of the Church, and in that he
hopes for the progressive enlightenment of the people (paulatim eruditur
vulgus), he expects also that an enlightened sovereign will some day
uphold the will of an enlightened people.

c. \textsc{Baruch} \textsc{Spinoza} (1632---1677) has in his \textit{Ethica more
geometrico demonstrata} (1677) made the most profound attempt
to work the fundamental thoughts of the new view of nature into a
general world-view. This work, despite its abstract
form, does not have an impersonal and purely theoretical character.
In Spinoza life and thinking go together. Clear thinking
was for him the way to spiritual freedom, to the highest form
of personal life. It was his endeavor to see all sides and forms
of existence in inner connection. The great thing
in his thinking is the full carrying through of different
lines of thought, each by itself, without interruption and without compromising,
and then the demonstration that each of the
trains of thought closed in themselves signifies only one side of or one
fundamental form of infinite existence. Thus he attempts
to show that unity and multiplicity, mind and matter,
eternity and time, value and reality each in its own
way express the whole of existence, so that there is an inner
identity between them.

Spinoza, who was born in Amsterdam of a Jewish merchant family
that had fled from Spain, began as a Jewish theologian, in whom
his co-religionists placed great hopes. But he took an
increasingly critical stance toward the religious notions of his fathers
and was at last solemnly expelled from the synagogue. He lived for
a time in the country, afterwards in the neighborhood of Leiden and finally in
% --- p. 27 ---
\setcounter{footnote}{0}%
\opage{27}The Hague, occupied with studies and with the working out of his
writings. His frugal livelihood he provided by grinding
optical glasses. A circle of young friends stood in lively exchange of thought
with him and studied his ``Ethics'' in manuscript.
His life is an exalted example of cheerful resignation and
inward devotion to intellectual work.

The deductive (``geometrical'') method that Spinoza employs
in his chief work is not the one by whose help he came
to his doctrine. From an unfinished treatise ``On the Improvement
of the Intellect,'' as well as from hints in the chief work,
one sees how he conceived the development of knowledge.
This begins with scattered and fortuitous experience
(experientia vaga), in which the various phenomena are put together
according to rules that present themselves involuntarily and rest
merely on the fact that no exceptions have been found. Science
(ratio), on the other hand, makes a careful comparison of the given
phenomena. It proceeds from experience, but seeks something
that expresses itself both in nature as a whole and in its individual parts:
general laws, valid everywhere. As examples Spinoza names
the laws of motion in the material domain and
the laws of the association of representations in the mental domain.
Only in such laws can one find a conclusion for one's
thoughts, whereas the series of individual phenomena can constantly
be continued, since what in one respect is cause, in
another respect is effect, and conversely. An Absolute in
which thought can rest is to be found only in the
law of the causal series, not in its supposed beginning or end. Such an
Absolute Spinoza calls substance: that which subsists in itself
and is understood through itself, in that its concept does not presuppose
any other concept. Spinoza's substance, the conclusion of his
thinking, is thus the principle of the lawfulness of existence.
The highest kind of knowledge is not knowledge of the laws of nature itself.
Above experientia vaga, which gives only individual cases,
and ratio, which gives general laws, there stands for Spinoza
intuitive knowledge (scientia intuitiva), by which the individual
phenomenon is immediately presented as a peculiar
member of the great connection of nature, --- the individual being
% --- p. 28 ---
\setcounter{footnote}{0}%
\opage{28}as part of the substance that reveals itself in this connection of nature.
This higher intuition can be attained only when
one has gone through the school of experience and of science.
It seems rather to be an artistic beholding than a
purely scientific conception. It gives a picture of the
individual seen in the light of the whole. Both in the second mode of knowledge
(ratio) and in the third (scientia intuitiva)
we regard things from the standpoint of eternity (sub specie
æterni), i.e. not in their isolation and contingency, but
as members of a comprehensive order.

Like Descartes and Hobbes, Spinoza builds his whole
philosophy on the causal principle, whose validity is for him, as
for them, self-evident. He lays quite special weight
on the point that cause and effect cannot be different things:
``When two things have nothing in common with one another, the one
cannot be the cause of the other; for there would then not be
anything in the effect that was also in the cause, and everything in the effect
would then have arisen from nothing.'' That two things stand
related as effect and cause means, according to Spinoza, that the
concept of the one can be derived from that of the other. He does not distinguish
between ground and cause. The relation between cause and effect
is for him the same as that between premises and conclusion
in an inference. That the cause precedes the effect \textit{in time,} not merely in thought,
he lays no weight on.

``From the standpoint of eternity'' time vanishes.

The cause of a thing can lie either in itself or in
something else. What has its cause in itself is substance
(fundamental being). It subsists in itself and can be understood through itself.
We have already seen that the fundamental being reveals itself for Spinoza
in the inner lawfulness of nature. It lies at the basis of
the subsistence and activity of all things. There can be only one
substance; for if there were several, they would have to limit one another,
and none of them could then be said to subsist in itself and to be understood
through itself. This concept, which for Spinoza forms the conclusion
of his thinking, is for him identical with the concept of God
as well as with the concept of nature. But these concepts
must then be understood in a way other than the usual one.
% --- p. 29 ---
\setcounter{footnote}{0}%
\opage{29}Nature signifies an inner power that stirs in everything, not the
mere sum of all that exists. And God cannot be conceived as
a personal being, since psychological properties such as intellect
and will cannot be applied to an absolute and infinite
being. Intellect presupposes an object different from it,
and will an end toward which one strives; but such
differences cannot hold for the absolute being. --- Things
that do not have their cause in themselves Spinoza calls \textit{modi}
(individual beings, phenomena). Modes have their cause in the substance,
whose forms of manifestation they are. They stand in external causal relations
to one another, but this whole series of causes and effects
given in experience has its cause in the substance, which is the inner bond, the
power that holds together the series of phenomena.

Everything that exists, modes as well as substance, appears
in our experience with two \textit{attributes} (fundamental properties
or fundamental forms): mind and matter (or, as Spinoza
expresses it, thinking and extension). Since substance is
an infinite and perfect being, it must have infinitely
many attributes; but we know only two of these, since experience
shows us no more. By attribute Spinoza understands ``that
which the intellect perceives of substance as constituting
its essence.'' Every attribute must therefore be understood
through itself, just like the substance whose essence it is to express.
One cannot therefore explain one attribute by another, mind
by matter or matter by mind. If one has begun to regard
existence under the attribute of matter, one must continue
with it and not leap over to another attribute. The essence
of existence (substance and modes) displays itself in the form of matter
as well as in the form of mind, and the one form cannot
be derived from the other. That follows already from the way
in which Spinoza conceives the causal principle: cause and
effect must be of the same kind, and so belong under the same
attribute. To derive a material change from the action of the soul
is therefore, according to Spinoza, in truth to admit that one
does not understand it. --- By this conception Spinoza founded
a new hypothesis, which leads beyond the opposition between
% --- p. 30 ---
\setcounter{footnote}{0}%
\opage{30}spiritualism and materialism, the identity hypothesis. Against Descartes
he objects that because two properties cannot be derived
the one from the other, it is not justified to refer
them to two different beings (substances); they may be
irreducible properties of the same being. From Hobbes he is distinguished
by the clear insight that the mental cannot
be derived from the material, and that the scientific
conception of existence therefore cannot be exhaustive.
--- In Descartes, Hobbes, and Spinoza we have the principal representatives
of the three principal hypotheses concerning the relation between mind
and matter. ---

Of Spinoza's special psychology his account
of feelings and drives is of lasting importance. Like Hobbes,
and under his influence, he starts from the impulse of self-preservation,
which in his system shows itself in connection with
the whole world-view, in that it is the infinite substance
that acts in the individual mode, whose striving for
self-preservation is thus a part of the divine activity.
When this striving is satisfied, pleasure arises; when it is
hindered, pain arises. Yet this happens only in the transition
to satisfaction or hindrance; a stable and unchanged
state would not be bound up with pleasure or
pain. The special kinds of feeling develop under the influence
of the association of representations, whereby a feeling toward
an object can gradually become attached to quite other
objects. Compassion can e.g. arise when we represent to ourselves
a being who resembles us as feeling pleasure or pain.

On the impulse of self-preservation Spinoza, like Hobbes, builds
his ethics. It is the opposition between activity and passivity,
between being cause and being effect, that here
becomes decisive. We are active when the cause of our state
and actions lies in ourselves, passive when it lies outside
us. As an individual being (mode) man is from the outset
predominantly determined by external conditions, by the relation to other
modes; his feelings and drives are called forth from without and make
him dependent. But by virtue of the impulse of self-preservation a striving
can be awakened to have, as far as possible, the cause
% --- p. 31 ---
\setcounter{footnote}{0}%
\opage{31}of his feelings and drives in himself. There then forms an
ideal of human life (idea hominis, tanqvam naturæ humanæ
exemplar), as it would be if it unfolded
in full freedom and independence. This goal furnishes a
standard for what is called good and evil. What leads
us nearer to the ideal is good; what removes us from it,
evil. These predicates can be applied only to finite beings
who stand in the midst of the struggle with external conditions. --- Since a
drive can be overcome only by another drive, the ideal itself must
awaken a drive that grows strong, so that it can come to
rule. Through life in society this drive is strengthened.
When men unite their powers, they can more easily
overcome external obstacles and provide for their physical self-preservation,
so that the mental goods, above all true knowledge,
can come to occupy the mind. The mental goods themselves
can be an object of common striving; they do not awaken the
strife that material goods do, which can each only fall to the share of
individuals; on the contrary, in the mental domain
the one can help the other without himself losing anything. When
full power and freedom are attained, the strong self-assertion
(fortitudo) in which virtue consists, then the individual will not only
with courage (animositas) assert his own personality,
but will also have power enough to give others, with magnanimity (generositas),
mental and material help. The highest degree of
freedom and self-assertion is, however, attained only through the full understanding
of ourselves in our unity with the substance of existence,
which stirs in us as well as in external nature. We then obtain a
scientia intuitiva of ourselves, see ourselves sub specie æterni, and
we are filled with a joy of knowledge that engenders love of
the Deity (amor intellectualis dei).

In his political theory Spinoza, like Hobbes, maintains a sharp
opposition between the state of nature and political life. But on
account of the significance that social life has, according to him, for
the development of the individual, it is so far from the case that the individual,
as in Hobbes, gives up his freedom by entering into
the state, that full freedom, especially in so far as it is conditioned by
mental development, first becomes possible in the state. The state is
% --- p. 32 ---
\setcounter{footnote}{0}%
\opage{32}not to make men into animals or machines, but to provide
conditions for the development of the mental and bodily
self-activity of men. It would then act against its own
end if it did not uphold freedom of thought, speech, and religion.
--- It is in his \textit{Tractatus theologico-politicus} (1670) that
Spinoza speaks for freedom of religion. In this work he maintains
at the same time the justification and significance of a purely historical
method in the study of the biblical writings

d. \textsc{Gottfried Wilhelm Leibniz} (1646---1716) is, like
his three predecessors, Descartes, Hobbes, and Spinoza, permeated
by the significance of the mechanical view of nature. As
a boy he had studied medieval works in the library of his father, a Leipzig professor,
and it was only after a
struggle that he joined the new natural science and philosophy.
He felt himself then ``transported into another world.'' He saw
that the new thoughts could not be rejected; but he could nevertheless
not give up the conviction that nature was in the end guided by
a Providence, and that the mechanical lawfulness was in the end
theologically grounded, i.e. had its ultimate ground in a divine
purpose. During a stay in Paris the physicist
Huyghens had great influence on his mathematical and physical
conception, and later he made personal acquaintance with
Spinoza. From 1676 on he lived in Hanover as electoral
court councillor and librarian, restlessly occupied with philosophical,
mathematical, historical, and juridical works. His universal
genius made it possible for him to intervene creatively and ordering
in very different domains. Everywhere he was guided by
one fundamental thought, continuity, the unbroken connection
through a multiplicity of transitions --- in the world of thought, of number,
of nature, and of purposes.

When Descartes and Spinoza taught that the sum of motion in
the world is always the same, Leibniz found in this the difficulty
that it was inexplicable how motion and
rest could succeed one another in the individual places in the world:
they stand, after all, as opposed states! Continuity can be
maintained only through the concept of force, i.e. that in a state which
% --- p. 33 ---
\setcounter{footnote}{0}%
\opage{33}contains the ground of a following state. It is force,
not motion, that is conserved. When motion ceases at
a place in the world, it is not succeeded by absolute rest, but by
a state of equilibrium, in which opposed forces hold one another
in equilibrium. The transition to motion then takes place by the
equilibrium being removed. Leibniz here laid the foundation of the
modern distinction between potential and actual energy, led by
his great interest in continuity, which in turn is connected
with the presupposition that everything has its sufficient reason.
If one were to ask again why force is conserved, Leibniz answers
that the opposite, a vanishing of force, would
be an imperfection and a decline, which conflicts with
Providence. He demands that mechanical explanation of nature be carried through
in the particular; but this last presupposition, the
conservation of force, is teleologically grounded.

On another point too Leibniz departs from his predecessors.
Descartes and Spinoza started from extension as
an absolute property (created by God or a divine attribute).
But the extended is always composite, and the truly
real must be the elements of which it is composed, and
which therefore (since they are to lie at the basis of extension)
cannot themselves be extended, but must be unities
(monads). If one wishes to form a representation of these monads,
no other way stands open than analogy with
our own souls. Every monad is a little world, within
which the following states have their ground in the preceding ones,
just as in the human soul mental states
develop out of one another through striving for clarity.
Leibniz takes the same stand toward the opposition between unconsciousness
and consciousness as toward the opposition between rest
and motion. He maintains continuity in existence, seen from
its mental side as well as from its material side.
Nothing in the world begins with absolute suddenness, by
a leap, consciousness no more than motion. Unconsciousness
is only a lower degree of mental life.

In his conception of the relation between the mental and the
material Leibniz stands on Spinoza's standpoint, in so far
% --- p. 34 ---
\setcounter{footnote}{0}%
\opage{34}as he embraces the identity hypothesis. But whereas for Spinoza
mind and matter stood as coordinate attributes, it is
according to Leibniz the mental side that is the fundamental one;
it expresses the innermost essence of existence, whereas matter
is its outer, sensible form. Material objects we must
interpret as sensible appearances of something that forms an
analogy to our own mental life. --- It is a significant
discovery that Leibniz has made, that all metaphysics, i.e. all conclusive
world-view, rests on analogy. All
mythologies, religions, and metaphysical systems had gone
this way, but Leibniz is the first to have become conscious of it. His own system is the first original attempt at
metaphysical idealism (i.e. the conception that the mental
constitutes the innermost kernel of existence) since Plato, and it stands
as a model for all later attempts in a similar direction. When
he himself has so great a confidence in analogy, it is because it
points out the only way to a conclusive understanding of the world,
and he proceeds, by virtue of the principle of sufficient reason, from the
assumption that the world must be intelligible. (\textit{Système nouveau} 1695.
\textit{Monadologie} 1714).

The law of continuity leads Leibniz to lay weight everywhere in nature and
in mental life on the small transitions, the vanishing
differences, in opposition to the leaps that experience
seems to show. The same thought guided him as a mathematician
and led him here to the invention of the differential calculus.
His ``differentials'' are vanishing small magnitudes
or changes of magnitude, which nevertheless by summation (integration)
give a finite magnitude. The differentials express
continuous variations. --- In many different domains
Leibniz's thinking has worked; but its type is everywhere
the same.

Among the infinitely many monads too there is continuity
with one another. They are all individually different. There are found,
Leibniz maintains, no two things in the world that are absolutely alike.
There would then (he adds teleologically) be no
reason why they should exist. In existence the greatest
possible multiplicity is realized together with the greatest
% --- p. 35 ---
\setcounter{footnote}{0}%
\opage{35}possible continuity. And the ultimate ground of this great
harmony Leibniz finds (in passing over from philosophy to
theology) in this, that all monads are emanations (``fulgurations'')
of a divine primal monad, which among all possibilities of worlds
has chosen the one that would exhibit the greatest possible
harmony. He even undertakes to prove (in his
\textit{Théodicée} 1710) that this world is the best of all possible
worlds. Every finite world must present imperfections,
partly limitations (metaphysical evils), partly
sufferings (physical evils), partly sins (moral evils). By
God's choice these imperfections are restricted to the
least possible. But beyond the domain of the possibilities even
the divine will cannot go. And when we find so much
evil in our experience, we must not forget that the part of the world
that we know, and that is perhaps the worst part, may
be vanishing in relation to infinity, --- and moreover,
that evil does not abolish the harmony of the world, but precisely
serves to bring it out, just as a dark shadow enhances
the effect of the bright colors in a painting, or as a sound
that alone would sound ill may fit excellently into
a whole piece of music! ---

Besides the works of Leibniz cited (only a few of the great
mass of notes that he has left behind) must be noted
\textit{Nouveaux Essais,} which appeared only long after his death,
and in which he countered John Locke's theory of knowledge.

\closerule

\chapterhead{III.}{THE PHILOSOPHY OF EXPERIENCE.}

The formation of systems in the seventeenth century had as its foundation
concepts and presuppositions that were not examined more closely. This holds
quite particularly of the causal principle
% --- p. 36 ---
\setcounter{footnote}{0}%
\opage{36}and of some propositions of natural science that, like it, came to be
regarded as self-evident. It is the merit of English philosophy
to have examined the psychological origin and the epistemological validity
of such concepts and presuppositions.
This inquiry acquired significance beyond the purely
philosophical domain, in that men now also put to the test the principles
that held in the state, in religion and in education. Authorities
that had hitherto won involuntary assent were now to render
an account of their origin and their justification.
Together with psychology and epistemology, ethics now
came to the fore.

a. \textsc{John Locke} (1632---1704) is said to have got the idea for his
epoch-making work \textit{Essay on human understanding} (1690)
from a discussion of religious and moral subjects; since
it proved difficult to reach agreement, Locke saw that one
must first examine our faculty of knowledge, in order to see what
it is capable of and what it is not capable of. He sets up the
principle that all our ideas are due to experience, to observation
either of external things or of our own inner states and
activities. This as regards origin; as regards validity,
he finds it self-evident that what we get from
observation, i.e. predominantly passively, must be valid, whereas
a difficulty arises with concepts and propositions
that are due to our own activity.

Locke was no more than Descartes and Hobbes satisfied with
his medieval instruction. The writings of the two men named
led him into philosophy. Since he could not subscribe to
the 39 Articles, he had to give up his original plan of
becoming a clergyman and afterwards studied medicine. Later he entered
the service of the Earl of Shaftesbury and remained attached
to his family, as tutor, secretary and friend. When the Earl was
overthrown, Locke went over to Holland, and he composed here
his most important writings. He returned to England with
William of Orange and gained no small influence under
the new government. He spent his last years in rural
solitude.

% --- p. 37 ---
\setcounter{footnote}{0}%
\opage{37}Since all ideas ultimately derive from experience, there
can be no question of ``innate'' ideas. The assumption
of such ideas is due to human laziness, which shuns the
trouble that may be bound up with discovering the origin of our ideas.
But from the simple ideas, which we receive
predominantly (for the most part) passively, we ourselves form new
ideas. We \textit{combine} simple ideas into complex ones;
thus we form ideas of things (beings,
substances) by combining the various qualities that
each for itself correspond to a sensation. We \textit{compare} different
ideas and thereby get new ideas, e.g. of
likeness and unlikeness, of causal relations. And we \textit{separate} (abstract)
ideas from their original connection, e.g.
when we form the idea of a color in general,
or of space apart from any content. How, then, can
such ideas produced by ourselves be
valid for existence itself? They are surely not copies or
impressions! Their validity can consist only in our being able to use them
as patterns or exemplars, by which we measure what is observed.
In this way we use in mathematics the geometrical
figures and in moral philosophy the ethical ideals. But now the concept
of thing (being, substance), which is expressly to concern
something that exists in itself? What we perceive is surely never
the things themselves, but only their qualities! It is no use
saying that there must be ``something'' underlying the qualities;
for then there must surely be ``something'' underlying this ``something'',
and so on. Nevertheless Locke will not deny the existence
of things or beings. --- Characteristic too is his attitude
toward the causal principle. He regards it as self-evident,
like the system-builders: Nothing cannot produce Something! And
he confidently uses it to prove the existence of God (since the world
must have a cause). The idea of cause, on the other hand, corresponds,
he thinks, to the observation that things and qualities arise
through the action of other things. But he has himself reckoned the idea of cause
among those that we ourselves produce (by comparison)!
--- Locke thus left problems enough to his successors.

% --- p. 38 ---
\setcounter{footnote}{0}%
\opage{38}b. \textsc{George} \textsc{Berkeley} (1685---1753) wished to remove the
consequences, dangerous to religion, of the new science,
and this end he thought to reach by a critique of the abstract
concepts (especially the concepts of space and matter) and by going back
to immediate experience and intuition. Seldom
have childlike piety and sharp critical analysis been united
in so close a manner as in this clear and fine mind.
At the University of Dublin he had studied Locke, Newton
and the chemist Boyle, and he wrote his principal works
as a young man. Then he became an Anglican clergyman and polemized
zealously against the freethinkers. His favorite plan was to
found a college in America, partly to convert the Indians,
partly to transplant art and science to a new flowering in
the new world. After three years' stay there he had to
give up this plan. He then served as Bishop of Cloyne
in Ireland, equally zealous as a pastor of souls, as a philanthropist and as a
patriot.

In his principal work \textit{(Principles of Knowledge.} 1710) Berkeley
shows that in the strict sense of the word we have no general ideas,
as Locke had still assumed. We
can represent to ourselves a part or a quality distinguished from other
parts and qualities. But we cannot represent to ourselves
anything ``common'' to several qualities, such as e.g. color in general.
An idea of color must be either an idea
of red or of yellow or of green, etc. When
I want to have an idea that is to hold of several qualitatively
different things, I must either regard a single one of these
things as an example, or use a word that fits each
of them. The idea of matter as something extended
and movable is thus a mere general idea; for we cannot
observe anything extended without color, hardness, etc.
And space (extension) we know (as Berkeley had
shown in \textit{Theory of Vision.} 1709) from two different senses, sight
and the sense of touch. Space in general (mathematical
space) is supposed to be ``common'' to sight and touch; but nothing
of the sort can be represented. Our idea of space must either
concern something that can be seen, or something that can be touched.

% --- p. 39 ---
\setcounter{footnote}{0}%
\opage{39}Berkeley wishes only to combat the assumption that behind
the sense qualities there should lie a material substance. He does
not wish to combat natural science itself, but sees its task
in finding out by what rules we, when one sense quality
has appeared, can expect another. To explain nature
is to point out the lawlike relations between the sense qualities.
--- --- As a theologian Berkeley declares it a detour
to let God first create ``matter'' and then let it
produce the sense-sensations in us. It is simplest
to think of God (whose action we represent to ourselves in analogy with
our willing, which constitutes the innermost essence of our soul)
as acting immediately in the production of our sense-sensations
according to definite laws.

c. \textsc{David Hume} (1711---1776) brought the empirical
inquiry into human knowledge to a provisional
conclusion, especially by his inquiry into the concept of cause, that
concept whose validity is a tacit presupposition for
research in natural science, psychology and history, as well
as for religious beliefs and for practical
human life. Hume's problems therefore lie at the root of
all human thinking. And he is a poser of problems of the first
rank, besides which his psychological talent makes it
possible for him to throw light on the actual development
of human thoughts even where their objective validity is problematic.

Hume was the son of a landowner in the south of Scotland and
showed early a desire and aptitude for thinking and learning. After vainly
trying to make his way in practical life
he withdrew into solitude and wrote, during a stay in
France, his principal work \textit{Treatise on Human Nature} (1739
---40). Later he turned to historical and economic
studies and wrote a history of England, one of the first works
in which account is taken of all sides of the development of culture.
In his \textit{Essays} (1748 ff.) there are, besides important
economic treatises, two philosophical treatises (on the
human understanding and on the principles of morals), in which
% --- p. 40 ---
\setcounter{footnote}{0}%
\opage{40}he gives shorter presentations of the most important problems of the principal
work. In \textit{Natural History of Religion} (1757) he treated
the religious problem from a psychological and historical standpoint;
in the \textit{Dialogues on Natural Religion}, written in 1751 but published only after his death,
he treated the same
problem from the critical side. --- After holding various
public offices he spent his last
years in Edinburgh in scholarly otium.

Whereas Locke had distinguished between the origin of ideas
and their validity, for Hume both coincide. An idea (idea) is valid only when it derives
from a sensation (impression). Whence the sensations
in turn come he does not inquire, because it is epistemologically
indifferent; with regard to validity it
comes to the same thing whether they derive from ourselves, or from external
things, or from God. Pure logic and mathematics also, which
concern only relations between our ideas, not the ideas
themselves, Hume excludes from his inquiry. He wishes
to discuss only the validity of those ideas by whose help
we think we can go beyond what is given in our
sensations. An idea of this kind is (as Berkeley
had already shown in regard to the material) the
idea of things (beings or substances). We sense
only single qualities, not that which ``has'' them. In
ourselves we observe changing sensations, ideas and
feelings, but no soul or ego, no constant
elements that could deserve the name of ``ourselves''. And
the concept of cause fares in a corresponding way. We observe
single phenomena that follow one another in time,
but an inner bond, a necessary connection we do
not observe. Causality cannot be observed. And all phenomena, all
things appear to our observation as independent, each with
its own constitution; we can very well represent them to ourselves each
for itself, without being compelled to pass from one to
the other. When a ball moves, we can represent to ourselves
this motion without being compelled to represent to ourselves any other
motion.

% --- p. 41 ---
\setcounter{footnote}{0}%
\opage{41}But we do after all use the concepts of thing and cause! Whence
can that come? --- Hume answers: The concepts of thing or
substance we form under the influence of the propensity that our
consciousness has to continue with an activity even after
the occasion has ceased. It is thereby that we form ideal
ideas: through our imagination we represent as
perfect what for observation is only imperfect.
Thus we form, e.g., ideas of absolutely regular
figures. And thus we form ideas of constant
things or substances, although we know only different degrees
of constancy. In the case of the concept of cause, custom
also takes a hand. We are accustomed to nothing happening without
something else having happened beforehand, and something third happening afterwards.
Therefore, when something happens, we expect partly a ``cause'', partly an
``effect''. But there is no proof in that. It is a mere association of ideas.
It goes with ``necessity'' as it goes
with the sense qualities; like these it is subjective, not
objective or real. Custom can ground a belief, an instinctive
conviction, but cannot yield proof. Hume naturally does not
mean that we could refrain from using the concept of cause;
but for him the marvel was that the validity of the
concept on which all our real science depends could not itself
be grounded.

Just as Hume in his epistemology ends by finding
the ultimate foundation in a need or an instinct whose justification
cannot be proved, so he also grounds his
ethics on feeling, not on reason. Reason can only compare
and analyze; it can convince us of facts
and relations, but not of value. Good and evil
are qualities that our feeling drives us to attribute to things,
just as our sense-sensations lead us to attribute to
them the sense qualities. Since we call actions and qualities
good even though they do not benefit ourselves, the
feeling that underlies approval cannot be of an
egoistic kind. In our approval or disapproval we take up a
common, not a private standpoint. When we, e.g., regard justice
as a good quality, this can only spring
% --- p. 42 ---
\setcounter{footnote}{0}%
\opage{42}from our embracing human life, wherever it stirs,
with sympathy. Perhaps it was originally only interest in
our own security that made us approve of justice; there
must then have taken place a change in our emotional life when we
call it good even where our own welfare is not affected by
its observance. The basic motive of ethical valuation is then
sympathy or the feeling of comradeship (fellow-feeling).

As in ethics, so in the philosophy of religion Hume opposes
the overvaluation of the influence of the intellectual.
And as in epistemology he undertakes here too
a twofold inquiry, in that he discusses partly the
psychological origin of religion, partly the validity of religious ideas.
The first is presented in ``The Natural History of Religion'',
the second in ``The Dialogues on Natural Religion''. --- Religion
arises not from purely intellectual motives, but from fear and
hope and from the propensity to represent all beings in
the likeness of man. In the beginning man represents to himself
the beings to which he takes refuge under the anguish of life, in very imperfect forms.
But here too the tendency of the imagination to go beyond the limited
and the finite is at work, and it dawns on man that his God
must be an infinite being, and that there can be only one
God. Alongside this tendency, which leads to representing
the Deity as raised above everything human and removed
from the finite world, there is, however, an opposite tendency,
a need to represent it as near and in a graphic
manner, in human form and with human qualities,
and we find in religion a constant oscillation between
these two tendencies. Thereby we understand the oppositions and
contradictions that are found partly between the different religions,
partly within the same religion. --- In his ``Dialogues'' Hume introduces
a supernaturalist (who maintains faith in a supernatural),
a rationalist (who thinks that by the way of reason
be able to prove the existence of God) and a naturalist (who maintains
that we can come beyond experience neither in faith nor in reason).
Although the naturalist in the end makes a polite retreat,
it is nevertheless clear that his grounds are decisive in Hume's
% --- p. 43 ---
\setcounter{footnote}{0}%
\opage{43}own eyes. He disputes the right to infer a God from
the order and purposiveness of nature. Why should what is
purposive in the world not be able to arise from natural
causes? Through a progressive development the
ordered wholes and systems in the world might be thought to have
formed themselves. And in any case the real world presents
so much that is imperfect that one cannot infer a
being who should be at once all-good and almighty.
Finally --- even if one wants to find the origin of the world in a
divine thought, one must ask again whence this
thought has its origin. --- When the naturalist nevertheless ends by
finding difficulties in all standpoints, we must presumably
not see in this mere caution, but an expression of Hume's
constant striving to pose the problems sharply and to keep
them open.

\closerule

\chapterhead{IV.}{THE PHILOSOPHY OF THE ENLIGHTENMENT.}

The systems and the philosophy of experience affected only a
narrow circle. But toward the middle of the eighteenth century
there awoke a striving to spread the ideas of the quiet thinkers
in larger circles. In England this movement, which
is usually called the Enlightenment, took on forms less characteristic
than in France and Germany. In France, where the
old order of things in church and state stood more closed and
unyielding in the face of all change than in England and Germany,
the movement of the Enlightenment took on a decidedly revolutionary character.
Locke's demand for a pointing out of the experiences
from which the ideas derive was used to call
the existing order to account, and men believed themselves to be
in possession of presuppositions sufficient for the
% --- p. 44 ---
\setcounter{footnote}{0}%
\opage{44}solution of the problems of thought and of life, so that whatever
did not agree with them had to be overthrown. In Germany the
after-effect of Leibniz's philosophy combined with Locke's
influence to produce a popular philosophy, which undertook
a purging of the traditional representations. But both
in France and in Germany there appeared men who were satisfied neither
by the new enlightenment nor by the old
tradition, because their gaze was directed toward more deeply lying
conditions of life. Such men were Rousseau in France and
Lessing in Germany.

1. On a visit to England \textsc{Voltaire} (1694---1778)
came to know English conditions and English thoughts. Through his
\textit{Lettres sur les Anglais} (1733, burned 1734) France came to know
Locke's and Newton's thoughts, as well as English religious,
literary and political conditions. Voltaire's poetic
activity and his tireless work for tolerance and
justice belong to the history of literature and of culture. Here only
his most important philosophical works shall be mentioned: \textit{Dictionnaire
philosophique portatif} (1764, burned the next year) and \textit{Le philosophe
ignorant} (1766). Now with arguments, now with mockery
and scorn, now with deep indignation he combated prejudices
and spiritual and bodily oppression. His arguments
he drew from Locke's doctrine of the origin of all ideas
from experience and from Newton's doctrine of the lawlike
order in nature, and criticized on their basis the doctrine of the
church. His own religious standpoint was the so-called natural
religion (cf. Herbert of Cherbury). He regarded that
religion as the best which contains much morality and few
dogmas. Reason can, from the purposiveness in nature,
prove the existence of God, but otherwise ends in insoluble problems;
we know neither what matter is nor what the soul
is. God has given us reason to guide us in our action,
not to penetrate into the essence of things. What lies beyond
natural religion is due to stupidity or deception;
stupidity makes the beginning, and deception then takes it into its
service. Educated people need no other religion than
the natural; the others had best go on going to Mass.
% --- p. 45 ---
\setcounter{footnote}{0}%
\opage{45}--- Although the purposiveness in nature proves for Voltaire
the existence of a God, there is nevertheless --- he thinks --- too much
evil and too much suffering in the world for this God to
be almighty. In his poem on the destruction of Lisbon and in the novel
``Candide'' Voltaire has expressed his thoughts about the
evil in the world. One may have hope of progress, but
the world is far from being ``the best world''. Voltaire
mocks Leibniz's optimism.

In decided opposition to Voltaire \textit{Holbach}
(1723---1789), a German baron living in Paris, presented materialism
as the true consequence of natural science. \textit{(Système de la
nature.} 1770). Motion is an original property of
matter, and material phenomena must be explained only by
material causes. It is then not necessary to assume
mental beings different from matter. Thinking or
consciousness is only motion in the smallest parts of matter. There is
only one science, physics, the doctrine of motion.
The assumption of a double kind of beings, mental and material, is
not merely useless, it is harmful. It favors superstition
and thereby the dominion of the priests. Even the so-called
natural religion is dangerous; for religion cannot subsist
without cult, and with that the dominion of the priests is at once
given. The representation of God is due to cunning priests (fabricateurs
de la divinité), who wanted by its help to rule over men!
---

Despite all opposition Voltaire and Holbach agreed
that the age of the Enlightenment had come, and that the only thing now
was to spread it as widely as possible. In Diderot's and
d'Alembert's great Encyclopedia, in which the most outstanding
men of the time were collaborators, the boldest thoughts
and the sharpest criticism found admittance under disguises and paraphrases.
After it the French philosophers of the Enlightenment
have got the name of Encyclopedists. To their circle belonged for a time
also \textsc{Jean Jacques Rousseau} (1712---1778), and his break
with them --- in which, besides the differences of principle,
many personal motives also played a part --- was a
cultural-historical event, a sign that a new problem
% --- p. 46 ---
\setcounter{footnote}{0}%
\opage{46}was about to break through. As Hume's problem concerned the possibility
of science, so Rousseau's problem concerned the value of
culture.

Rousseau was born in Geneva. His restless spirit, yearning beyond the
usual limits of life, cast him into a
romantic life of wandering and adventure, which in the year 1741 led him
to Paris, where he became a friend of Diderot and Holbach. When
the great thought of the opposition between nature and culture
had dawned on him, he left Paris to live in the country,
and soon afterward broke with his encyclopedic friends.
On account of his writings (which were burned like Voltaire's,
but were not anonymous like his) he became an exile;
even in Switzerland he could find no lasting place. In
his later years his suspiciousness and fear of persecution
had definitely the character of insanity. He spent his last years
in France in seclusion.

By nature Rousseau understands the immediate and unified
force of life, the involuntary unfolding, in opposition to the
divisiveness and limitation that culture causes, especially through
the division of labor, which does not let the individual attain all-round
development, but fosters a single faculty at the expense of the others,
whereby men then become more dependent
on one another. The healthy, involuntary self-assertion (amour
de soi) does not stand in opposition to compassion and devotion;
on the contrary, one in whom the inner life stirs
vigorously needs to make no distinction between himself and the others. It
is only when an unhealthy comparison and rivalry assert
themselves that self-assertion becomes egoism (amour propre).
To return to nature means for Rousseau to return
to the unity and fullness of life and of the forces, where neither
inner nor outer oppositions assert themselves. In opposition
to the divisive effect of culture he exclaims: Rendez
l'homme un! In free nature, which Rousseau had really learned
to know in the wanderings of his youth, this feeling of unity
stirred especially vividly in him. He is the first writer
who has really given expression to the joy in the solitude of nature
and the beauty of nature. --- Since feeling --- in opposition
% --- p. 47 ---
\setcounter{footnote}{0}%
\opage{47}to knowledge and will, thus to comparing and preferring
--- has quite particularly the character of immediacy,
unity and fullness, he saw man's innermost essence and worth
expressed in it, and it is due especially to him that it is now regarded
as an independent side of mental life.

The irreconcilable opposition between nature and culture, which
Rousseau maintained in his first writings (especially \textit{Discours sur
l'origine et les fondements de l'inégalité parmi les hommes,}
1755), was, as his thoughts deepened and
became clearer, replaced by a relative opposition. In his later writings
\textit{(Émile} and \textit{Contrat social,} 1762) he distinguishes between a
natural and an unnatural culture. Culture consists in the development
of the forces, and a culture is natural when it corresponds to the stage of development
at which man stands, so that self-activity
becomes possible and the unified force of life does not come to suffer
from the division of labor. What Rousseau protested against
was that what the age praised as culture should signify an
unconditional progress, and next that culture should without more ado
be capable of being transferred from one people to another. In the
pedagogical, the religious and the social domain alike he maintains
the involuntary unfolding of life as the one thing needful.
--- All influence from without is to serve only to draw out the
distinctive character of each individual; education is to get the child
to develop his faculties on his own, and has its chief significance
in keeping obstacles away. Childhood is an
independent period of life and must not be treated merely as
a means or preparation for the adult age. --- In religious
respect Rousseau agrees with Voltaire, as far as
the representations are concerned. He believes in a good God, who is, however,
not almighty, since the evil in the world can be explained only
by the resistance that matter offers to the divine action.
Yet he takes a stand not only against materialism, but
also against Voltaire. He felt a great antipathy to Voltaire's
standpoint, which had a purely intellectual character. For
him himself religion had its root in the life of feeling. The need
of self-unfolding can be so great that it seeks beyond
every limit and finds rest only in an infinite object.
% --- p. 48 ---
\setcounter{footnote}{0}%
\opage{48}The religious feeling reaches its culmination in an ecstasy of
love, admiration and overflowing devotion. That religion
springs from nature means for Rousseau that it
springs from inner feeling. In a letter he says: J'ai
laissé là la raison, et j'ai consulté la nature, c'est-à-dire le
sentiment intérieur qui dirige ma croyance. --- In his political theory
he maintains the sovereignty of the people. The general will (la
volonté générale, different from la volonté de tous) is to rule.
It is the innermost need of the people and its dominant tendency, directed
toward the common interests, the welfare of the whole and of the individuals
through the changing generations. To obey it is not dependence,
for each individual is membre du souverain and can
find his will fulfilled in it.

2. In Germany it was especially \textsc{Chr. Wolff's} (1679---1754)
popularization of Leibniz's philosophy that came to lie at the
foundation of the Enlightenment. His broad, systematic writings
displaced medieval philosophy and theology from the German
universities. It was not, as with Leibniz, metaphysical
idealism and the doctrine of monads that stood in the foreground, but the doctrine
of the harmony in the world effected by God. His firm faith
that there was a sufficient reason for everything made him a
martyr, in that King Friedrich Wilhelm I of Prussia, on account
of an accusation of fatalism, deposed him from his professorship
at Halle and --- with the threat of the gallows --- ordered him to leave
the country within 48 hours. Wolff then went to Marburg, whence
Friedrich II called him back immediately after his accession to the throne.
--- Alongside Wolff's popularization of the systems of the seventeenth century
Locke's influence also worked in Germany.
German philosophy of the Enlightenment became then an often very external
combination of Wolff and Locke. Gradually the
psychological interest decidedly came into the foreground before the metaphysical,
and men were thereby led to see that the opposition between
the obscurity and the clarity of representations was not the only
decisive one. Under the influence of English literature and
Rousseau, feeling began to be recognized as an independent
side of mental life. The shadow side of this was the sentimentality
that marked the age alongside rationalism.

% --- p. 49 ---
\setcounter{footnote}{0}%
\opage{49}In the strife between rationalism and orthodoxy that
resulted from the appearance of the philosophy of the Enlightenment in Germany,
\textsc{Gotthold} \textsc{Ephraim} \textsc{Lessing} (1729---1781) took up a position
that pointed beyond the oppositions. Although in character
and gifts he is very different from Rousseau, he
nevertheless occupies an analogous position in the history of the life of the spirit. He
had a great and fine sense for what is original in the domain of the spirit,
the religious and the poetic as well as the intellectual,
and for an inner connection that lies behind the
definitely marked forms of faith and thinking. Therefore
he took up a critical position toward both rationalists and orthodox.
As a theological critic he pointed to original
Christianity, which was older than the so much disputed
Bible (\textit{Über den Beweis des Geistes und der Kraft}). And he
set, as an investigator, the constant striving after truth, in which
the forces come to living unfolding, above the
mere possession of truth, which makes one lazy and passive. \textit{(Duplik).}
The connection in spiritual development he found not in
dogmas or results, but in the inner striving to which they owe
their origin. The sense for what is original and distinctive led
him also as an aesthetician. He referred to Shakespeare as the
great exemplar of dramatic art, in opposition to the prevailing
French taste, which had ended in the artificial and empty.

Lessing's own religious standpoint is expressed in the proposition, directed against the
orthodox, that from historical reports one
cannot ground a knowledge of the eternal connection of existence.
The various positive religions stood for him
as stages in the spiritual development of man; what education is
for the individual, the positive religions are for the species. The
Old and the New Testament are ``the elementary books of the human race''.
The time will come when such books will be
unnecessary; then ``the third kingdom'', the new eternal gospel,
will come. \textit{(Erziehung des Menschengeschlechts.} --- \textit{Gespräche
über die Freimaurer.)} --- In his general philosophical
ideas he stood close to Spinoza and declared that if he were to
name himself after anyone, he would name himself after him.

% --- p. 50 ---
\setcounter{footnote}{0}%
\closerule

\opage{50}\chapterhead{V.}{IMMANUEL KANT AND THE CRITICAL PHILOSOPHY.}

Like Lessing and Rousseau, Kant too belonged to the century of the ``Enlightenment,''
but like them he pointed beyond it by the standpoints
he adopted. Like so many great minds
he received influences from many sides.
Wolff, Newton, Hume and Rousseau had an especial influence
on his philosophy, as he presented it, after a long period
of thinking for himself, in his principal work \textit{Kritik
der reinen Vernunft} (1781). The task of this work is an examination
of the powers of human thought to attain science
without relying on experience; by ``pure reason'' is meant
precisely thought, as it works without the aid
of experience. That all knowledge \textit{begins with} experience, Kant
does not doubt; but it did not follow from that that it
\textit{need not be founded on} experience. The philosophers of the seventeenth century
had assumed that there could be science beyond the domain
of experience, but on the other hand Hume
had maintained that the very principle that underlies all science
of existence, the causal principle, rested only on
a custom and could not be scientifically grounded. It was
precisely Hume who, as Kant himself has put it, awakened him
from his dogmatic slumber. For in his first period, as
an adherent of Wolff's philosophy, he had trustfully relied on the fundamental
propositions of science and on the justification of applying
them beyond experience. The influence of Hume introduced
for Kant the long labor that ended with the definitive
founding of the critical philosophy, whose principal task
is an investigation of the presuppositions of human science
and of their validity. Kant is aware that the
critical philosophy cannot be man's first philosophy.
From the first, man does not doubt the
% --- p. 51 ---
\setcounter{footnote}{0}%
\opage{51}validity of his knowledge and does not assume that it has limits. The first philosophy
is therefore always \textit{dogmatic,} i.e. it makes no
examination of the fundamental concepts and fundamental propositions from which
it starts, but relies on them trustfully and believes that with their
help it can solve all riddles. When the dogmatists come
into conflict with one another, a \textit{skeptical} period naturally follows,
in which it is doubted whether valid
knowledge can be had at all. But the \textit{critical} philosophy then sets itself the task
of judging between dogmatists and skeptics by a methodical
examination of the conditions of science and an investigation of
the domains in which these conditions are present. --- It is not,
as the title of Kant's principal work shows, all science
that he wishes to investigate; that which rests merely on experience, e.g.
descriptive natural science and psychology, he does
not include.

Besides Wolff and Hume, Newton
and Rousseau, as already said, exercised a profound influence on Kant's development.
--- \textsc{Isaac} \textsc{Newton} (1642---1727) stood for him as
the model of the strict scientific conception of
the world. In his \textit{Principia philosophiæ naturalis mathematica}
(1687) Newton had given the greatest example of exact science
combining induction and deduction. Starting
from the experience that gravity is greater in the valleys than on
mountain heights, and that bodies thrown into the air fall down,
he conjectured that the planets too are heavy, and that they deviate
from the direction in which they would move according to the principle of inertia
only as much as the law of falling bodies established by Galileo
on the surface of the Earth indicates. He then drew the mathematical consequences
of this conjecture and showed finally that the results
of this deduction agreed with the observations.
The planets thus move according to the same law as bodies
thrown into the air from the Earth. And from this assumption could
then be derived the laws found by Kepler on the basis of Tycho Brahe's observations
(among them the law of the elliptical
form of the planetary orbits). It now became Kant's task to find what
presuppositions Newton's science rested on, and to investigate
whether these presuppositions could be proved. But he
% --- p. 52 ---
\setcounter{footnote}{0}%
\opage{52}also tried to continue Newton's work. The latter had
declared that the wonderful harmony in our solar system ---
in that planets and satellites move in the same direction, roughly
in the same plane and in almost circular orbits --- could
not be explained by natural causes. Kant showed (in Allgemeine
\textit{Naturgeschichte und Theorie des Himmels} 1755) that if
one imagined that all the globes of the solar system had once
formed one great rotating nebular mass, one would be able to explain
the order now existing according to ordinary laws of nature.
During the rotation, by reason of centrifugal force,
rings would separate off from the nebula's equator; these would in turn split
into globes at various distances, but all rotating in the same
direction and in roughly the same plane. It was the same hypothesis
that (independently of Kant) was later set up by Laplace,
but which is now disputed. It was a principal thought of
Kant that all phenomena are subject to the law of causality, and that even
the order and purposiveness found in nature must
be explicable according to laws of nature. All the greater a problem arose
for him when, later, through Hume's criticism, he came
to doubt that the causal principle could be proved at all.
--- Rousseau's influence led Kant to a new valuation of man:
formerly he had regarded knowledge and enlightenment
as the highest and had despised the rabble, who know nothing;
the study of Rousseau brought him to find an inner kernel
in man's feeling and will, which can subsist with a certain
independence of enlightenment and information. Rousseau also exercised
an influence on the later development of Kant's ethics.
---

Kant's life (1724---1804) moved, as far as externals are concerned,
in a narrow circle. His father was a poor craftsman
in Königsberg. At the University of Königsberg he studied
Wolff's philosophy and Newton's physics. After having for
some years been a private tutor in various noble families
in East Prussia, he worked for a number of years as Privatdozent
at the University of Königsberg with great following;
only in 1770 did he become professor. He was never outside
East Prussia; his time was taken up with the preparation
% --- p. 53 ---
\setcounter{footnote}{0}%
\opage{53}of his works and of his academic teaching. Yet
he took part in social life and had a reputation as a witty
and elegant man of society. He belonged to the age of the Enlightenment, but
by ``enlightenment'' understood a striving, not something that one could
attain once and for all, still less had now attained. And while
his critical thought led him into depths that ordinary
enlightenment did not know, he had a sense of the sublime,
which was the keynote of his view of life, nourished
by consideration of nature and by his conviction of the unconditional validity
of the ethical ideals. In his old age he came to
suffer under the clerical reaction after Friedrich II's death.
When in the year 1793 he had published a work on the philosophy of religion,
a royal rescript was addressed to him, which censured his
free conception and threatened further measures
if he continued in the same direction. Kant considered it his
duty to obey, but declared that he would now not write
or speak about religious subjects at all. When a new, more liberal government came, he resumed his activity in the philosophy of religion
and made public the scandalous episode. His last
years present the picture of the dissolution of a great mind; probably
he suffered from a brain disease; only occasional glimpses
still recalled what he had been. ---

Kant's philosophy can be presented in three sections, according to his three
most important works, of which the first gives an epistemology,
the second an ethics, and the third seeks to reach beyond the
results of both in a conclusive conception.

1. \textit{Epistemology.} (Kritik der reinen Vernunft).

a. Since Kant wishes to investigate how far we can get along
the path of pure reason, he must first find out what pure
reason is. He then analyzes our knowledge and finds that there
is something that constantly asserts itself, whatever we know or
experience, and which we ourselves must bring with us, namely the forms
in which we combine and order our observations. Knowledge is
a mental activity whose laws and forms of operation must
be traceable in all experience and all knowledge. In contrast
to these forms Kant calls the mere sensations and observations
the material of our knowledge. That we receive passively.
% --- p. 54 ---
\setcounter{footnote}{0}%
\opage{54}It is only the forms that Kant, in accordance with his task, has to do
with, and of them he finds that three kinds in fact assert
themselves, which all, however, present the same fundamental character and
are manifestations of one and the same fundamental form.

In every sense perception there is something that is independent of the
particular sense qualities I perceive: in one and the same
picture of space I can combine differently colored things at different
times, and in the same length of time I can experience highly different
events. Everything I perceive, however different
and changing it may be, I combine in space or in time
or in both forms. Space and time are \textit{forms of intuition,}
by whose application the chaos of sensations is combined
and ordered into sense images.

But we do not remain with the sense images. These
are in turn combined and ordered by reflection (``the understanding'')
and placed in definite relations to one another. Here too
certain forms assert themselves: the forms of the understanding, or \textit{categories}
(fundamental concepts). By a very artificial path, which we
cannot enter into here, Kant seeks to show that there are 12
such categories; but he indicates that they can be reduced
to two: the category of magnitude and the category of cause. \textit{The concept of magnitude}
we apply when we proceed from the assumption that everything we experience must
have a certain extension in space and a certain degree of intensity, and
that between all points in space and between all degrees
on the scale of intensity a continuous transition is possible.
\textit{The concept of cause} we apply when we proceed from the assumption that the events
we observe must be capable of being connected in such a way that there
is not merely a temporal relation between them, but also such a
relation that the one event contains for us
the ground of the other, so that there is an analogy between
the relation ground---consequent on the one hand and the temporal relation
between the events on the other. Both categories
express the demand for \textit{continuous connection} and
work together to make our experience able to constitute a unity. By
virtue of the concept of magnitude we pass continuously from place to
place and from degree to degree, and by virtue of the concept of cause we pass
continuously from event to event. There shall be
% --- p. 55 ---
\setcounter{footnote}{0}%
\opage{55}no leaps in nature, neither in extension nor in
degree nor in the series of causes.

Both forms of intuition and categories are forms of
our knowing activity of consciousness, which at every stage
shows itself as a combining, a \textit{synthesis.} They are different
ways of combining the given material. No experience
is possible except in that the sensations are connected into sense images,
and the sense images are connected into continuous series
in time, space, degree and causality. Such synthesis need not
take place with full consciousness. It is, says Kant, the effect
of a blind but indispensable mental function, without which
we should have no knowledge at all, but of which we only
seldom become conscious. All analysis we can undertake presupposes
that something is connected in our consciousness; for we
can only dissolve what is connected.

Forms of intuition and categories express connections of
the given material. But the involuntary need, deeply
grounded in consciousness, to connect and combine is so strong
that it will embrace everything in a concluding whole. A conclusion
of all extensive and intensive series and of the series of causes
cannot be thought within experience, which is always limited.
Insofar as synthesis leads to concepts that require completion
by the formation of absolute wholes, it forms \textit{ideas,}
a third group of forms. Such an idea is the concept
``world,'' which is to embrace everything existing, although no limit
can be set for space, time and the series of causes. Likewise
the concept ``God'' in the sense of a first cause that does not
in turn have any cause, and the concept ``soul'' in the sense of a
substance that lies behind consciousness and cannot be discovered
through self-observation. It is in and of itself natural that such
ideas are formed; our striving to combine
is one with our innermost nature.

b. The investigation by which the forms of the activity of knowing
are found is descriptive. It can only show
in what way our knowledge in fact works, but
not the validity and justification of the way in which it works.
Kant then makes the transition from psychology to% --- p. 56 ---
\setcounter{footnote}{0}%
\opage{56}epistemology (``critique of reason'') by asking how a knowledge
that rests merely on the forms, not on the material, --- thus a
knowledge of ``pure reason,'' can be science. The problem
arises for Kant (as already for Locke) from the fact that knowledge
is an activity on our part; with what right
do we ascribe to the forms in which this activity appears
validity for things? --- Kant's answer is that if it turns out
that the application of our forms of knowledge is the \textit{condition
for our being able to have real experience} (as opposed to mere
representation, fantasy or dream), then propositions that
are built on these forms must also hold for \textit{all possible experience}
--- but not any further. He therefore investigates in this respect
the various groups of forms.

All content of experience is apprehended \textit{in space or time or both
in space and time.} Since now pure mathematics properly develops only
the laws of space and time, it must hold for everything that
we can ever experience, for all possible observations. From this
follows the validity of applied mathematics, notwithstanding that mathematics
in itself is a pure science of reason. We can know in
advance (a priori) that all experiences must agree with
mathematical laws. But since space and time are human
forms of apprehension, we dare not say that the knowledge we here
reach is the only possible one. We know existence only from
the standpoint of man. What Copernicus maintained, that in
apprehending the motions of the heavenly bodies we must not
forget that we see them from the human stage, holds
for all our knowledge, all our science, because it is due to
a labor that is carried out in human forms.

The case is similar with the \textit{categories.}
Since they are forms without which we do not understand things, everything
that we are to experience must fall under them. Every phenomenon
or event must satisfy the demands of the concept of magnitude:
their arising and development must proceed
continuously, so that a higher degree of intensity is always
reached by traversing all lower degrees of intensity, and a
more distant point in space is always reached by traversing
all points that lie between the starting point and it. If
% --- p. 57 ---
\setcounter{footnote}{0}%
\opage{57}such a continuous transition cannot be shown,
we cannot distinguish between experience and imagination, reality
and dream. The same holds for the concept of cause.
If I cannot arrange the events in series, within which each member
is related to the following in an analogous way as, in our thinking, ground is related
to consequent, I cannot have real experience. When
I can imagine the events in any order whatever,
arbitrarily make them change places, I cannot be certain
whether I am confronted with something real. In my consideration,
in my memory and my fantasy the representations can come
in different order; but when my representations are to express
a reality given in experience, they must form a
series in which the members cannot be interchanged. In my consideration
or my fantasy I can just as well let the house
come into being for me from above downward as from below upward; but if
the coming into being of the house is to be a real event, only one
single order of succession is possible. (Cf. Psykologi V D, 2). Everything that
we shall ever be able to experience must be subject to the causal principle.
Although this proposition is a proposition of pure reason,
in that it rests on one of the forms of our combining
activity, it must thus be declared valid in advance (a priori)
for all experiences. --- --- But the validity of the categories is,
like that of the forms of intuition, limited to possible experience.
The grounding goes no further. Experience is, says Kant, the
synthesis that gives validity to every other synthesis.
In the case of the ideas the grounding which in the case of forms of intuition
and categories led to the goal cannot be applied. For the
peculiarity of the ideas was that they demanded a concluding
synthesis, an unconditioned whole, a first or last. But experience
cannot present such a thing, since experience is always
limited. The ideas cannot therefore be said to specify conditions
for distinguishing between experience and imagination. The world,
e.g., cannot appear in any experience; beyond every
whole we form by synthesis, a more comprehensive whole can
always be thought. On the ideas, therefore, no
science can be built. There is no science of God, of
% --- p. 58 ---
\setcounter{footnote}{0}%
\opage{58}the world as an absolute whole or of the soul as substance.
Kant recognizes an empirical psychology; but it is no pure
science of reason. --- The significance of the ideas lies, on the other hand, in this, that
they point beyond every single experience, prevent us from settling
down content with any result whatever; they set ever new tasks.
Here Kant meets Lessing in the thought of the constant
striving after truth as opposed to the mere possession.
---

When Kant grounds the real significance of the knowledge of reason
on this, that it expresses the conditions for having experience,
he takes the concept of experience in the strict sense as a firm and
necessary connection of observations. He has not seen that
experience in this sense is itself an ``idea'' or an ideal,
since observations satisfy the demands of the forms of intuition
and of the categories only imperfectly and approximately. Kant
has therefore furnished no conclusive proof of the validity of the
causal principle and has not solved Hume's problem.

The limitation he showed for pure reason in
its various forms he derived from this, that we have no right
to regard these forms as the only ones in which existence
could be apprehended and understood. We must distinguish between the
way in which we know and the essence of things themselves. We
know things only as they appear to us, according to our
forms, not as they are in themselves, or as Kant expresses
it, we know only phenomena (Erscheinungen) not
the thing in itself (das Ding an sich). This he grounds
also by this, that whereas the forms are our ways of working
in our knowledge, we receive the material (the sense qualities)
purely passively. Whereas the forms belong to our nature, the material
cannot be derived from it. The concept ``Ding an sich'' then serves
to mark a limit. It is a \textit{limiting concept,} which establishes
firmly, partly that our forms are \textit{not} the only possible ones, partly that
the material of our knowledge \textit{cannot} be derived from its form. It is
thus a negative concept. But involuntarily Kant comes
to ascribe to it a positive significance, in that he calls ``Ding an
sich'' the cause of phenomena (especially as regards the material, yet
also as regards our forms). Hereby he comes into
% --- p. 59 ---
\setcounter{footnote}{0}%
\opage{59}contradiction with his own doctrine that one has no right to apply
the concept of cause beyond experience; for ``Ding an sich''
cannot be experienced. And besides, the question arises
whether we can be sure that ``Ding an sich'' will always
give us material that fits our forms, and will always let
the same forms subsist in us. On this point too,
not merely in the concept of experience, it appears that Kant has not
freed himself from all dogmatism.

2. \textit{Ethics.} (Kritik der praktischen Vernunft. 1788).

Under strong influence from Rousseau, as well as from
Hume and other English philosophers, Kant in his first
utterances on ethical questions laid much weight on the feeling of
``the beauty and dignity of human nature'' as
the foundation of the ethical principles \textit{(Betrachtungen} über das
Schöne \textit{und Erhabene} 1764). Later he laid more weight
on the will, self-activity and on the necessity that
in our striving we be always in agreement with ourselves.
His ethics thereby acquired a predominantly formal character. When now
historical studies led him to regard the striving of individuals
as members of the striving of the whole species to
raise itself from the mere state of nature to true culture,
the ethical principle could not merely become the agreement of the individual
with himself, but had to become one that all
individuals in the same situation could consistently follow. I shall
then act in such a way that the principle of my action can be set up
as a rule for all. I shall, e.g., return goods entrusted to me,
even if I could keep them without being discovered;
for the opposite principle cannot be set up as a rule, since then
no one would entrust anything to others. Every human being is to be treated
as an end, never merely as a means; only in that
way does a human society become possible, in which the
freedom of the individual can subsist with the freedom of others. Just as we
obtain valid knowledge only when all sense-sensations (all
``material'') are ordered according to laws that ultimately spring from the
forms of our mind, so our willing and acting are valid only when
they do not follow feelings and drives (whether these are
sympathetic or egoistic), but submit to a
% --- p. 60 ---
\setcounter{footnote}{0}%
\opage{60}law of reason that can hold for all. Such a law makes itself known as
an unconditional command (a categorical imperative) in face of all other
incentives, in face of all authorities, and in face of everything that
experience teaches or seems to teach. Most clearly of all
the ethical appears in such cases where duty and inclination
stand in the greatest opposition to one another. There will always,
Kant thinks, be an element of the feeling of pain in the consciousness of duty
owing to the resistance which the sensuous parts of our
nature offer to subordinating themselves to the form of the law.

In the capacity thus, in spite of all inner and outer resistance, to
subordinate one's acting to a universally valid law Kant sees
the mark of nobility of man. The law is not, after all, an alien law,
but grounded in man's nature as a rational being, one with
his innermost will, however exalted its authority may be.
Through the ethical law man is a citizen of another
world than that of experience. Inwardness (in its springing from man's
innermost nature) and exaltedness (in its opposition to
the human nature appearing in experience) are the two
characteristic features of the ethical law according to Kant's conception.
He takes this law to be purely formal and a priori, sprung
from pure reason. But it clearly presupposes
that a human society is to exist, and properly formulates
only its most important conditions. It was also
in Kant himself historical considerations that had led him
to his final ethical conception. ---

According to Kant there is an inner relation between religion and
ethics. The transition from the ethical to the religious he finds
in this, that man is to maintain the good, that which the unconditional
ethical law demands, under inner and outer resistance.
From this follows a sharp opposition between duty and happiness,
between ideal and reality, freedom and nature. This opposition
will in turn give rise to a need to believe in another existence,
in which the striving that is broken off in the world of experience
can be continued in ever higher forms, and to believe in
a being who is able to bring duty and happiness, ideal and
reality into harmony. The belief in God and in immortality,
the content of so-called natural religion, cannot, however, be
% --- p. 61 ---
\setcounter{footnote}{0}%
\opage{61}scientifically grounded. It is practical postulates that lead to
it, postulates which cannot be refuted, since science
itself has its limit (in ``Ding an sich''). Their significance
consists in their keeping us up in our striving to maintain
the validity of the ethical law. But they do not extend our knowledge;
that follows already from this, that the forms of our knowledge
hold only within experience. All religious representations
are therefore grounded only on analogy; they are symbolic and
figurative. --- In his work \textit{Religion innerhalb der Grenzen der
blossen Vernunft} (1793), --- the work that brought Kant into
conflict with the authority of the state, --- he has carried through a
symbolic conception of the Christian dogmas, in which he
finds deep ethical truths contained.

3. \textit{Theory of Unity on an Aesthetic and Biological Basis.} (Kritik
der Urtheilskraft. 1790).

In his epistemology and in his ethics Kant had maintained
a sharp opposition between the world of experience and ``the things
in themselves.'' Indeed his whole philosophy is characterized by sharp
distinctions, which in his opinion were necessary if
the validity of knowledge and the unconditional character of the ethical ideal
were to be maintained. But precisely according to the principles of the critical
philosophy, which indeed impress the illegitimacy of transferring
our thoughts to existence in itself, the
thought had to arise whether those very distinctions and
oppositions might not have to be charged to the account of our human
way of apprehending. Kant found two groups of facts
that could point in the direction of a unity, the one group
of an aesthetic, the other of a biological kind.

a. In the phenomena we call beautiful and sublime,
existence produces in us a disinterested satisfaction. It is
not our sensuousness that is gratified when we stand before
something beautiful. The beautiful rests on this, that our faculty of intuition
and our understanding can easily and involuntarily connect the given
multiplicity into a unity, just as if the phenomenon were adapted
to our faculty of apprehension. Kant here thinks especially of
beauty of form, as it appears in harmony and symmetry,
in the works of nature and of art. In the case of the sublime our
faculty of apprehension and our self-feeling are
% --- p. 62 ---
\setcounter{footnote}{0}%
\opage{62}overwhelmed, in that we stand
before what is immeasurable in extension and force; but precisely
thereby is awakened the consciousness of a force dwelling in ourselves, striving beyond
all limits, that which expresses itself in the ethical
law. Both in the beautiful and in the sublime it appears
thus that there is something in existence that agrees with that
in us which is raised above sensuousness and egoism. --- In genius,
in genial production, there is a feature that points in
the same direction. For in genius an involuntary force of nature
and natural impulse is at work, and yet works are thereby produced that stand
as models. In genius nature works as if it were
conscious art. Here too, then, the opposition between
man and existence does not assert itself.

b. In organic life we have a fact that can be
interpreted in an analogous way to the beautiful, the sublime and
genius. Nature works in the organic in a way for
which we properly have no expression. We see here a
whole that does not arise by external joining of the parts,
and in which, scientifically regarded, we also have no right
to assume a preceding plan or a purpose (as in the
coming into being of a human building). Nature works here too, as if it were
conscious art. Neither the teleological
nor the mechanical explanation of nature can interpret the organism.
But the opposition between teleology and mechanics perhaps rests
only on the nature of our understanding: a whole we can only
understand when either the parts are joined together by external causes,
or a plan precedes. In existence itself, on the other hand, the
same order of nature that makes the causal connection
possible could also make possible the arising of organisms, purposively
constructed forms of life. In the unknown ground of nature
teleology and mechanism could be one. It was a thought on which Kant
had already in his youth, when he worked out his hypothesis on the
development of the solar system, laid so great weight.

And in the same way it might be with the
greatest opposition: between the laws of nature that science shows,
and the ideals that the ethical law leads us to recognize.
It would then be one and the same principle that made itself known
% --- p. 63 ---
\setcounter{footnote}{0}%
\opage{63}in the laws of nature and in the ethical ideals. It would then not merely
be through postulates that that opposition could be overcome.

Kant is, however, aware that such a theory of unity cannot
be carried through scientifically.

\closerule

\chapterhead{VI.}{THE PHILOSOPHY OF ROMANTICISM.}

a. Modern philosophy, as we have seen, received its task and
its character determined partly by the new world-picture and the
new science, partly by the freer conception of human
relations that came with the Renaissance. Later,
Kant's philosophy was determined not only by the opposition between Wolff and
Hume, but also by Newton's natural science and Rousseau's
problem of culture. The development of German philosophy in the
period immediately after Kant was determined by the influence of the
aesthetic and religious tendency that in the history of literature is called
Romanticism, and which marks a strong reaction against the
Enlightenment and Revolution of the eighteenth century. People were tired
of criticism and analysis and felt the need to immerse themselves immediately
in nature and in history, to feel one with the innermost forces
of existence. The split between poetry, religion, and
science, this fruit of the division of labor in the spiritual
domain, was to cease; the spiritual life was to constitute a
unity, and poetry, religion, and science were to be only
different forms of this unity. People wanted, in general,
to go beyond all oppositions, and therefore also beyond Kant's
philosophy, which had emphasized the oppositions so strongly.
They took their starting point from the thought that Kant had so cautiously
suggested, of a unity behind the oppositions at which we stop,
--- a thought that for Kant meant the limit of thought. And
% --- p. 64 ---
\setcounter{footnote}{0}%
\opage{64}it then became the ideal to derive everything from one principle. In their enthusiasm
for this ideal people did not examine its justification
or the possibility of realizing it, and overlooked that every
inference, after all, must be built on at least two premises.
The methods of thinking and of science that had hitherto been valid
they despised, and demanded a speculative thinking that could
solve the highest problems. And they broke the connection
with natural science that philosophy had had since the days of Bruno
and Descartes. Starting from intimations of a unity in everything, they now wished
themselves to construct the forms of nature as a series of
steps up to the world of spirit. What they arrived at by this route
was sometimes poetry, often fantasy, not science. And
yet the philosophy of Romanticism has not only uttered profound thoughts
of lasting value, but also paved the way for an understanding
of the connection among the forces of nature, which
could later be carried through in a more exact manner, just as
it also contributed to awakening and nourishing historical sense.

The three most important representatives of the philosophy of Romanticism
are Fichte, Schelling, and Hegel.

J. G. \textsc{Fichte} (1762---1814), a Saxon peasant's son, followed in
his early youth with lively interest the appearance of Lessing
and later, after heroically working his way through
straitened years of study, was led to philosophy by the study of Kant's writings.
At the University of Jena he worked as a teacher with
great success, not only because of his power of thought and
eloquence, but through his whole energetic personality. When
he was dismissed on account of his religious views, he went
to Berlin, where he became one of the first teachers of the new university.
He was one of the men who worked for Prussia's
national regeneration after the battle of Jena, especially through his
``Reden an die deutsche Nation'' delivered during the French occupation.

Fichte took his starting point from Kant's opposition between
thinking and ``Ding an sich''. Thinking is an activity,
an action, and things are partly objects of this
action, partly its products. But how then is it to be
explained that things exist as something different from the activity of our
% --- p. 65 ---
\setcounter{footnote}{0}%
\opage{65}spirit? Things as objects cannot be derived
from the thinking and acting ego itself; it cannot
itself call forth its limits, its resistance. The only solution
is that things (objects, non-ego) exist only in order that
we may have something to act upon, something that can release
our forces. The free use of the forces is the highest thing;
but it is possible only through there being limits, resistance to
overcome. --- --- It is thus for Fichte an ethical consideration that
gives the solution. Nature is for him material for our action in accordance with duty,
our self-development to true freedom.

\textsc{Fr}. W. \textsc{Schelling} (1775---1854) is the real
representative of Romanticism in philosophy. Without any critical scruple
he proclaimed in his youthful writings, which are the
most significant he has produced, a new science that was to
solve all riddles. He too worked for a time in Jena, which was
the gathering place of Romanticism, later at various Bavarian
universities. In his old age he was called to Berlin by the Romantic
King Friedrich Wilhelm IV to counteract
the radical Hegelians; but his appearance there was a
disappointment of great expectations.

Schelling made an impression on his time at the beginning of the nineteenth century
through his ``philosophy of nature'', by which he understood
a spirited beholding for which nature stands as unconscious
spirit, which by various steps through the inorganic and the
organic world approaches the point where consciousness
awakens to progressive clarity. We can know
nature only because it is spirit itself at a lower stage.
Nature must be understood from within, not, as ordinary natural science
attempts, through external interaction between
the different matters and forces. --- Yet it was no
real development that Schelling assumed in nature. The different stages
exist simultaneously; it is only for thought that their kinship
shows itself. And it was a symbolic interpretation, not a scientific
explanation, that was given; the principle vera causa was
pushed aside. The enthusiasm that this ``philosophy of nature'' aroused
became in part important for the development of poetry (as with
% --- p. 66 ---
\setcounter{footnote}{0}%
\opage{66}us through the lectures that Schelling's disciple gave in Copenhagen
1902); but for thinking it was a regress.

\textsc{G. W. F. Hegel} (1770---1831), the systematizer of Romanticism,
also worked for a time in Jena. After the battle of Jena he went
to Bavaria and later came to the University of Berlin, where he
founded a numerous and influential philosophical school.

Neither Fichte's ethical postulate nor Schelling's beholding
satisfied Hegel. He set himself the task of developing
the thoughts that express the innermost nature of existence, in such a way
that they form a series in which each member leads with necessity
over to another. No single thought can exhaust
existence; every thought shows itself, when consistently carried through,
to be limited, but precisely by being driven to
the extreme leads over into another thought; if these two
thoughts are now combined, we get a higher unity, of which they are both members.
But this new thought has itself its opposite, together
with which it forms a new higher unity, and so
on --- a progressive process, always through triads.
This is what Hegel calls the dialectical method.
But it is not only in our subjective thinking that dialectic
appears. There is also dialectic in existence. Development
proceeds through oppositions and the overcoming of oppositions.
Thus the organic germ unfolds in interaction
with the environment, and in the fully developed organism
we have the higher unity of what was laid down in the germ
and of the opposition to the environment. The immediacy of childhood
reaches, through the doubt and strife of youth, the
maturity of manhood.

Hegel's system itself consists of three members. In pure logic
the inner relations of thoughts are presented. In the philosophy of nature
external relations in space and time take the place of the pure relations of thought.
Through this externality thought develops into
the philosophy of spirit, in which the opposition between the inner and
the outer falls away. In art, science, and religion,
the highest forms of the life of the spirit, existence reaches its consummation.
And it is through all these stages and forms the absolute
world-spirit that unfolds itself.

% --- p. 67 ---
\setcounter{footnote}{0}%
\opage{67}In Hegel's ethics the idea of the state is the main thought. His ethics is
a social ethics. The good has its subsistence in society, in
which forces are at work that lie beyond the consciousness
of the individual. The individual reaches his highest development only by
living in and for society --- in the family and civil
society, but especially in the state, which takes up both into itself. The state
cannot be explained by an association of individuals, but is
a higher order of things, in relation to which the single
individual is of vanishing value. Against all individual
ideals and utopias Hegel emphasizes that the ethical idea is not
so impotent that it should be unable to be realized; it
has found its real form in the historically developed state. It
is not the business of philosophy to construct ideals, but to demonstrate
the ideality in the results of historical development. --- In
an analogous manner, according to Hegel, philosophy is related to religion.
There is only a difference of form between philosophy and
religion; the content is the same. Philosophy teaches the unity
in existence through all oppositions and at all
stages; religion teaches that everything has its origin in God. The former
speaks the language of thought, the latter that of the image. What religion proclaims
in great images and in a concrete manner, philosophy expresses
in the form of thought. Especially in the doctrine of the Trinity and the doctrine
of the Incarnation Hegel found the same truths
expressed that he had developed dialectically in his system: the Trinity
expresses that the highest cannot be an abstract
unity, but a unity of oppositions; the Incarnation
testifies that the Deity, precisely by passing
over into its opposite, into humiliation and suffering, shows itself
as the highest world-power. It was especially in theology that
Hegel's philosophy was received as an ally; one
saw in it a form in which the theological doctrinal content
could receive a new grounding. As Hegel thought he had reconciled
all oppositions, he thought also that he had reconciled
faith and knowledge.

b. The Romantic tendencies of thought prevailed in Germany and
the neighboring countries until toward the middle of the century, and they still
continue to assert themselves in altered forms. It is the
% --- p. 68 ---
\setcounter{footnote}{0}%
\opage{68}interest in asserting an idealist conception of existence
that still leads many to see in them a high point in the history
of philosophy. Kant's critical line of thought did not, however,
disappear; in the very period of Romanticism it held its own as an
undercurrent, represented by able investigators, such as \textsc{Fries},
\textsc{Herbart}, and \textsc{Beneke}, whose significance lies especially in the
domain of psychology. Besides, there were men who, though
strongly moved by the tendencies of Romanticism, nevertheless
maintained an epistemology in Kant's spirit, while they sought, by
different paths, a view of life that could lead beyond
the limits of science. Such men were, besides the
theologian \textsc{Schleiermacher}, \textsc{Schopenhauer} in Germany and
\textsc{Kierkegaard} in Denmark. Finally, there were men who, proceeding
from Hegel's school, but in sharp contrast to the Master's conservatism
in the philosophy of law and the philosophy of religion, developed radical
views. David Strauss maintained that on
Hegel's principles one could not distinguish between theology and
mythology. \textsc{Karl} \textsc{Marx} and \textsc{Ferdinand} Lassalle developed
Hegel's political theory in a socialist direction. \textsc{Feuerbach} led the philosophy of religion
back to pure psychology and suggested a philosophical
conception that is related to the positivism that arose in France,
which we shall bring forward in the next section.

Only a few of these thinkers can we dwell on here:
Schopenhauer, Kierkegaard, and Feuerbach.

α. \textsc{Arthur} \textsc{Schopenhauer} (1788---1860) was the son of a rich
merchant in Danzig, received a many-sided education, and learned early
to know the world through travels and by moving among
people. Independent as he was, he could devote himself entirely
to his studies and to the development of his view of life.
After an unsuccessful attempt as an academic teacher in Berlin
he lived as a private man in Frankfurt am Main. Very early
he had, through his own passionate nature, come to know
the dark and disharmonious forces and drives of life,
and what he saw around him aroused partly his indignation, partly
his pity. Above these experiences, which taught him that
not wonder (as Plato thought), but consternation and
% --- p. 69 ---
\setcounter{footnote}{0}%
\opage{69}sorrow was the beginning of philosophy, he sought to raise himself
by intellectual work and by artistic beholding.

Schopenhauer protests against the view that Fichte, Schelling, and
Hegel were Kant's rightful successors, and that the critical
philosophy was a stage left behind. In the first book of
his chief work \textit{(Die Welt als Wille und Vorstellung,} first
part 1819, second volume 1844) he presents an epistemology
that is closely related to Kant's. He is then clear that
knowledge (die Vorstellung) cannot teach us to know the innermost of
existence (das Ding an sich); but he thinks he has
found another way to the solution of the riddle. We ourselves
belong, after all, to existence; it stirs in our own inner being.
And now there stirs, far deeper down than all knowledge
(Vorstellung), a yearning and craving, an impulse of self-preservation,
that makes itself known in all our pleasure and pain, in hope and fear, in
love and hate, --- --- a will that constitutes our innermost essence,
the primal phenomenon in us. All representations are from the first
servants of this will. What is the primal phenomenon in us must
--- thinks Schopenhauer --- also be the primal phenomenon in existence
as a whole; by the innermost in us we understand the innermost
in the world. Thus he makes the transition from epistemology
to metaphysics by means of an analogy, the justification of which he
in his Romantic mood does not doubt. --- Against the
weight that Hegel laid on the life of thought, it was of great importance
that the life of the will was emphasized in so energetic, and in particulars
often so brilliant, a manner. A voluntarism was here set
against intellectualism. Already Fichte, by whom Schopenhauer
is influenced, had, by the way, given the will a predominant
place in his philosophy.

If it is asked what the will is aiming at, Schopenhauer answers:
at existing for the sake of existence itself; it is a will
to life (Wille zum Leben). There is here no longer any
Why. The principle of sufficient reason holds only
for knowledge (Vorstellung), not for the will. The manifold
forms and forces of nature in their mutual clash and
struggle, the ever renewed processes of motion, and the constant
unrest in the world bear witness to the
% --- p. 70 ---
\setcounter{footnote}{0}%
\opage{70}never-resting energy of the blind impulse of self-preservation.
This dark need beguiles us with the
illusion that life is good and valuable. But life appears to
us as good only because the will to life rules in us at any price.
In itself existence is pain, precisely because its innermost
essence consists in a never-resting and never-sated
need, and only as a phenomenon of contrast does pleasure or
satisfaction arise, namely when the subterranean fire is for a time
damped. All pleasure is illusion, is a zero point, which only because
of the contrast to the constant ardor and unrest comes
to stand as something positive (Cf. on this Psykologi VI E, 6).
In lively descriptions of human and animal life
Schopenhauer depicts the torments of existence, the mutual
struggle of living beings. These torments the great mass does not notice,
or easily forgets again. Those who see more deeply, especially the brilliant
human beings, on the other hand, experience the disharmonies and tear asunder
the veil of illusion.

The will to life can, however, be abolished by three different paths.
One can place oneself in relation to life as a mere spectator, by
being entirely absorbed in aesthetic or intellectual beholding. When we
e.g. immerse ourselves in the contemplation of a work of art, our
will is calmed, and we forget that we belong to the world. It is
chiefly this path that Schopenhauer himself has taken. But also in
love of humanity (which because of the torment of life must
appear as pity) the will is abolished, in that I make
myself entirely one with others. On this thought Schopenhauer
has built his ethics. But the complete abolition of the will
succeeds only for the saint, the ascetic, for whom there
is no longer any motive for action. In Buddhism, in
primitive Christianity, and in mysticism Schopenhauer finds the
highest practical solution of the riddle of life and the torment of life, and
he stands with deep admiration before the great representatives of asceticism
--- all the more as he was aware
that he himself was no saint.

β. In Denmark Hegel's philosophy had gained no small
influence, especially among aestheticians and theologians (J. L. \textsc{Heiberg},
and H. L. \textsc{Martensen}). F. C. \textsc{Sibbern} (1785---1872) gave an
apt criticism of it and himself put forward a world-view
% --- p. 71 ---
\setcounter{footnote}{0}%
\opage{71}that laid great weight on the fact that development in nature and in human life
proceeds from scattered points (sporadically) and for the
most part through the struggle and clash that arise when
different tendencies of development meet --- in the inner life of man
or in nature and in society. Since every human being is himself
one of the sporadic starting points, he cannot survey existence;
hence the debate between different standpoints in
the history of thought. In particulars Sibbern has especially occupied
himself with psychology, for which his power of observation and
his fresh sympathy with human life made him
well suited.

The religious thinker \textsc{Søren} \textsc{Kierkegaard} (1813---1855)
renewed the critical philosophy and disputed Hegel's solution of
the problem of knowledge. We can have system, i.e. order
and coherence in our thoughts, but we cannot comprehend
existence in a closed system of thought. We understand only
what lies ready and completed before us; but we live in the midst
of time, and our thought always comes after the experience, since
new experiences always come. We live forwards, but we
understand backwards. While we think, life goes on, in us and
outside us. And what holds of scientific thinking holds
even more of the thinking that takes part in our personal
view of life. Here the point is, first and foremost, that the thoughts
be our own, acquired or appropriated at first
hand; only then do they correspond entirely to our personality and have
personal truth. He who prays with inwardness to an
idol prays nevertheless to the true God, whereas he who in indifference
and out of habit prays to the true God prays to an
idol. Personality (subjectivity) is the truth. Kierkegaard
sympathizes with Lessing's words about the constant striving
after truth. --- It is in his chief philosophical work
Uvidenskabelig \textit{Efterskrift} (1846) that Kierkegaard develops his
conception of knowledge and sets up the concept of personal truth.
Here he also presents his comparative philosophy of life,
the doctrine of the different ``stages'' or types of personal
view of life and conduct of life. Three of these, the aesthetic, the
ethical, and the religious, he has depicted in poetic form in
% --- p. 72 ---
\setcounter{footnote}{0}%
\opage{72}\textit{Stadier paa Livets Vej} (1845) (in ``Enten-Eller'' he depicts
the first two of these stages). The standard for the ``stages''
Kierkegaard takes from the degree of the oppositions that the
spiritual life is able to embrace. The aesthetic man lives in the pleasures
of the moment; the ethical man has an inner core that he
faithfully preserves amid the changes of life, and he finds his task
in society (family, state); the religious man lives in the consciousness
of the cutting opposition between eternity
and life in time. The transition from one of these stages to
another cannot, according to Kierkegaard, take place by continuous development,
but only by a leap, which cannot be psychologically
motivated. (Cf. on the transition from the aesthetic to the
ethical Psykologi VI E, 4). --- More and more, for Kierkegaard
it was not the energy of the synthesis itself, but the suffering that the
strong tension between the oppositions caused, that became the real
standard, and he thereby came into opposition to Romanticism's
proclamation of the harmony of all oppositions, as well
as to the usual conception of Christianity.
Thereby was called forth the violent struggle against the established church
that filled his last years, and which he finally conducted through
the pamphlet \textit{Øjeblikket} (1855).

Kierkegaard lived as a private man, absorbed in the great
philosophical, aesthetic, and religious production, in which his
sharp reflection and rich emotional life, his melancholy and his
passionateness found an outlet. When he had said in ``Øjeblikket''
his last word, he succumbed to illness.

γ. \textsc{Ludwig} \textsc{Feuerbach} (1804---1872) had, under Hegel's influence,
passed over from theology to philosophy. But he could
not remain with Hegel's philosophy of religion, which sought to
find philosophical ideas in the religious representations. He set
himself instead the task of finding the psychological source of
these representations, of giving a psychological explanation of the origin of dogmas.
He resumes Hume's psychology of religion.
His chief work is \textit{Das Wesen des Christenthums} (1841).
--- Feuerbach lived, after having tried a career at the
University of Erlangen, as a private man in rural solitude,
% --- p. 73 ---
\setcounter{footnote}{0}%
\opage{73}occupied with studies and authorship. In his last
years he had to struggle with poverty and illness.

Religions arise, according to Feuerbach, in that the inward
yearning for the fulfillment of the wishes of the heart bursts the limits
of reason. Thereby it is explained that religion at its
height acquires an antirational character. Wish is the theogonic
(god-creating) principle. The religious predicates are
the content of human wishes; men's representations
of the nature of heaven and of the qualities of the gods
are expressions of what moves in the innermost of their hearts.
That God is personal means that the personal life is of
highest value (is divine); that God is love means
that love is of highest value (is divine).
If one wants to understand religion, one must make the religious predicates
into subjects.

From the beginning man has no reason to
set limits to his wishes and to the representations that the wishes
call forth. And it lies in the nature of feeling to make its
object infinite and yet to hold fast to its reality. Doubt
arises only when man comes to know his limits,
and when he begins to distinguish between his representations and
reality.

Between religion and culture there exists an inverse relation.
They are two different paths by which man hopes
to reach his purposes, and the more confidence he has in the one,
the less use he has for the other. In a similar manner
it stands with religion and ethics. The more the difference
between God and man is impressed, the more must qualities
such as love and justice be taken in a quite
different meaning when they are applied to God and when they are applied
to men, and a sharp conflict can arise
between God's will and what human love and justice
demand.

No real values disappear with religion. The
qualities that are valuable can continue to be asserted,
even if they are not ascribed to gods. Feuerbach grounds ethics
on love of humanity, in which the single individual's
% --- p. 74 ---
\setcounter{footnote}{0}%
\opage{74}unity with the whole species makes itself known. He lays weight on the fact that
the individual must know happiness and unhappiness from his own experience
in order to be able to have compassion with others, and therefore asserts
the justification of the impulse toward happiness. --- In his general conception of
philosophy (expressed in \textit{Grundsätze der Philosophie der Zukunft}
1843) he finds its object in man as a being of experience
and in nature as the foundation of his existence. Against
materialism he emphasizes (like Schopenhauer) that
matter is a concept formed by man, and that therefore the
conception that regards man as a product of matter
cannot be conclusive.

\closerule

\chapterhead{VII.}{POSITIVISM.}

The two great tendencies of thought in the nineteenth century
are Romanticism and positivism. The former proceeds from the forms
and ideals of thought, the latter from the factually given; ``positive'' means
above all ``real, posited, given.'' However different, indeed opposed,
they are, they both mean, each in its own way, an
opposition to the century of the Enlightenment, of criticism, and of revolution.
The positivists, like the Romantics in their own way, wanted
to immerse themselves in nature and in history. Both tendencies
sprang up at the beginning of the century, and positivism
is not to be understood as a reaction against Romanticism, although
it came to the foreground only when the dominion of Romanticism
was broken. --- Whereas the philosophy of Romanticism had its chief seat
in Germany, positivism belonged essentially in
France and England.

1. Positivism has its origin in the social school
founded by \textsc{Saint Simon}
(1760---1825). His task was to
prepare a social reformation, which in his view
% --- p. 75 ---
\setcounter{footnote}{0}%
\opage{75}was possible only if one founded a new world-view,
which for the present could be the same as Christianity
had been for the Middle Ages. Such a new world-view
had to be built on the foundation of the positive sciences,
just as men too could exploit nature in common labor only on this foundation,
instead of, as before, having
exploited one another. For this it was required that the theological
presuppositions that had formerly been laid at the foundation of science
be eliminated, so that it was built on purely natural presuppositions.
When this had happened, the positive philosophy could be founded.

\textsc{Auguste Comte} (1798---1857), the real founder
of positivism, studied at the Polytechnic School in Paris;
but when this was closed by the Bourbons on account of the students' revolutionary
disposition, he continued his studies
on his own and gave them a more encyclopedic character,
to which his collaboration with Saint-Simon also contributed.
This collaboration ended because Saint-Simon, in Comte's
opinion, subordinated science too much to his ideas
of social reformation and the religion of the future. With great
capacity for work and power of concentration Comte carried through his
presentation of the positive philosophy as a great encyclopedia
of the positive sciences \textit{(Cours de Philosophie positive.}
1830---1842). In his later years his thinking took on a more
subjective and mystical character, in that he regarded himself as
the founder of a religion of humanity and even instituted a
kind of cult.

According to Comte human knowledge (as a whole
and in the individual domains) passes through three stages. At the theological
stage only a small circle of experiences is at one's disposal,
and imagination therefore plays a great role. The
bond that here unites the facts into a whole is the notion
of spirits and gods. The cultural-historical significance of theology
rests on the fact that at this stage of development it has served
science as an intellectual bond, just as it also in the
religious authorities furnished a foundation for morality. At the
metaphysical stage the explanation is found no longer in
% --- p. 76 ---
\setcounter{footnote}{0}%
\opage{76}personal beings, but in general forces or ideas. One speaks
of force of attraction, chemical force, vital force, etc. This
is a progress, since such forces or ideas are assumed to
act in a regular manner. But this stage has
for Comte a predominantly negative character. It is the stage of individualism
and of revolution. At the positive stage
imagination and speculation are subordinated to experience.
Now the laws of phenomena are sought, i.e. the constant
relations between them. The different laws cannot be traced
back to a single one; the unity of science consists only in
its being the same chief method that is applied in all domains.
This unity in method makes possible a spiritual community
such as had not been known since the Middle Ages.

The chief method appears, however, in somewhat different manner
in the different sciences. These can be arranged in a series,
so that one begins with those that concern the simplest
subjects and set up the most comprehensive laws; these are also
the ones that have soonest reached the positive stage, and in
which deduction can be applied to a great extent. The further
we advance in the series, the more complicated the subjects become,
the more special the laws become, the more must induction
be applied, and the later the positive stage is reached. The order
is: mathematics, astronomy, physics, chemistry, biology, and sociology
(theory of society). There is no continuous transition
from one science to the other. --- --- We miss in this
classification logic, which would have its place before mathematics,
and psychology, which in Comte is divided between biology
and sociology.

Comte distinguishes between the positive facts and their
connection. This connection takes place at the positive
stage partly by ordering according to likeness (par similitude), partly
by deriving the following from the preceding (par
filiation), but in both cases it is a matter of a synthesizing
(tout se reduit toujours à lier). (\textit{Discours sur l'esprit positif.}
1844). Comte therefore also counted Kant among his predecessors.
--- In his ethics he takes up a similar standpoint to Grotius
and Hume. The concept of duty springs from the individual's
% --- p. 77 ---
\setcounter{footnote}{0}%
\opage{77}consciousness of being a member of a social whole, a consciousness
that arises when the solidarity of all human life is experienced and recognized.
Involuntarily man follows the need for intercourse,
and the utility of this dawns on him only afterward.
The germ of the feeling of solidarity lies already in the relation between the two
sexes and in care for the offspring; in the human world
there then goes on a constant education toward altruism (a
word coined by Comte to denote the opposite of egoism).
The single, isolated individual is only an abstraction. ---

2. In a wider sense the main tendency of nineteenth-century
English philosophy too is positivistic. At the middle
of the century \textsc{John} \textsc{Stuart} \textsc{Mill} (1806---1873) took a
highly esteemed place in the scientific world through his works
on logic \textit{(System of Logic} 1843), his political economy,
and his ethical and social-reformatory works. He
was strongly influenced by the radical ethicist and social reformer
\textsc{Jeremy} \textsc{Bentham}, but also by Comte, as well as by
the Romantic tendency, which in England was represented by
Coleridge and Carlyle. Mill held an office in the service of the East India
Company and used his leisure for his studies.
In his interesting autobiography he has given an account of his
course of development, in which so many different tendencies
intervened, without his quite being able to work them together.

In his logic he takes up a sharper opposition to Kant
than Comte did. He tries to conceive all thinking as
mere association, especially as association by contiguity. The basic form
of all inference he finds in the fact that when two
phenomena have frequently occurred together, we expect the second
when the first has occurred. It is an inference from particular
to particular. When many such experiences have been made, we formulate
them into general propositions, which can form the
foundation for deduction. But all deduction presupposes an
induction, ultimately a transition from particular to particular.
The causal principle too is grounded on induction. --- --- Mill's
logic is the counterpart of Hegel's. Whereas Hegel wanted to derive
all thought-content from the forms of thought, Mill wanted to derive the
forms of thought from the given content. Here the
% --- p. 78 ---
\setcounter{footnote}{0}%
\opage{78}opposition between Romanticism and positivism appears in the sharpest
way. --- The significance of Mill's logic lies in his characterization
of the special methods and in the honest discussion of
the questions of principle.

As an ethicist Mill, following Bentham, maintained ``the morality of utility''
(utilitarianism), which values actions according to
the degree in which they work toward the greatest possible happiness for
as great a number of people as possible (a principle already set up by Bentham
and earlier philosophers). The ethical feeling
that drives one to act according to this principle develops
under the influence of various motives, among which sympathy
occupies an eminent place.

3. A great extension positivism received through the hypothesis of evolution,
which appeared shortly after the middle of the century. Already
before Darwin's publication of his theory Herbert
Spencer had maintained that in considering the experiences to which human
thoughts and feelings owe their evolution, one must not think merely
of the experiences of the single individual, one must take
account of the species from which it descends. But it was especially
through Darwin's connection of the doctrine of evolution with the idea of
the struggle for life that this point of view won more general
attention.

\textsc{Charles} \textsc{Darwin} (1809---1881) receives with right a place
in the history of philosophy, partly because his results (like
Copernicus's, Galilei's, and Newton's) are of such great importance
for the world-view, partly because his conception of science
and its limits has philosophical interest. Darwin himself
lived as a quiet investigator in rural solitude, after he in
his youth on a three-year voyage around the world had collected
a great stock of observations. --- His explanation of
the origin of species \textit{(Origin of Species} 1859) is quite in the spirit
of positivism. To assume that each species is created by God is
for Darwin only to say that one does not know how it has
come into being. Darwin rests instead on the fact that organisms
vary, and that these variations can be more or less
beneficial to the organism when it is to procure food,
light, and warmth and to secure its offspring and itself against dangers and
% --- p. 79 ---
\setcounter{footnote}{0}%
\opage{79}enemies, or, as Darwin with a figurative expression calls
it, when it is to hold its own in the struggle for existence. The variations
that were decidedly unfavorable in this struggle would gradually
die out; only the more or less purposive ones
could hold their own. This doctrine had for Darwin the significance
of serving as a thread of thought (the intelligible thread), by which
a great multitude of facts could be connected. A limit
for his hypothesis he saw in the question of the origin of
the variations between which the struggle for existence (or,
as Darwin also calls it, natural selection) stands.
A similar riddle he finds in the origin of life in general.
--- Later investigators have maintained that only through considerable
``leaping'' variations (mutations, as their discoverer
Hugo de Vries has called them) do new species come into being, in the case of which
it then becomes a question whether they can hold their own in the struggle
for existence; this is thus not properly the cause of the origin
of new species.

From Darwin's autobiography and letters \textit{(Life and Letters of
Charles Darwin} 1887) one sees that his last standpoint with
regard to the relation between teleology and mechanism was
similar to Kant's. Just as little as he could believe that the world,
as we know it, was the work of a conscious purpose
(design), just as little could he rest content with its
being due to chance or blind force (chance, brute
force). He declared himself an agnostic (i.e. one
who acknowledges his ignorance concerning that which lies beyond
experience).

\textsc{Herbert} \textsc{Spencer} (1820---1903) passed from practical life
over to philosophical inquiry. He began as an engineer, then
got an interest in social questions and was thereby led
to psychological and philosophical studies. His philosophy sprang
from the need to find a fundamental thought that could synthesize
a series of natural-scientific, psychological, and sociological
studies. In his \textit{Essays} he has given a series of in part
brilliantly written treatises on individual questions.
Later he published in ten volumes his ``synthetic philosophy,'' whose
% --- p. 80 ---
\setcounter{footnote}{0}%
\opage{80}first volume \textit{(First Principles} 1861) gives the foundation for his
world-view: the concept of evolution.

In his theory of knowledge Spencer shows, under the influence of
the Scottish Kantian William Hamilton, that only that which
stands in relation to something else can be known. An unconditioned
and all-embracing thing can therefore not be known. Every piece of knowledge
gives only a section of an infinite reality, which can only
be expressed by means of symbols, not in the form of the concept.
Spencer is, on these grounds, an agnostic like Darwin.
--- But all the knowledge we can gain gathers
around the concept of evolution. We have understood a phenomenon only
when we have followed it through its different stages to its
present shape. It turns out that here there are certain features
that constantly recur, no matter which phenomena we
have before us, and which together serve to define what we
understand by the word evolution. There are three such main characteristics.

1) Evolution is concentration. A transition takes place from
looser to firmer connection when a phenomenon comes into being.
Thus when a heap of sand forms on the shore,
in the coming into being of star systems and globes, in the evolution of the life of
consciousness and of social life. --- 2) But next,
evolution is differentiation, transition from a more uniform to
a more diverse condition. Whereas concentration takes place
by uniform parts being acted on by the same force
(the grains of sand e.g. by the wind), differentiation takes place by
the uniform parts within the whole formed by concentration
being acted on by different forces (as when a
part lies in shade, another in sun). Within the whole there then comes
to be difference between the parts. Within the solar system
the individual globes take shape, and on the same globe there comes to be
difference between the inner and the outer parts, and between the
different parts of the surface. In the organism greater
differences appear between the organs the further growth
advances, and organic life on earth appears at
later stages with more different species than at earlier stages.
So it goes also with the life of consciousness, where the
% --- p. 81 ---
\setcounter{footnote}{0}%
\opage{81}differences at later stages assert themselves most strongly, and in society
the division of labor is an expression of the increasing differentiation.
--- 3) When concentration does not fall away through
differentiation, but the firm connection is preserved despite
the arising of differences, one passes over from a more indefinite
to a definitely ordered condition. Here evolution appears
as determination.

This third form of evolution reminds one of what Hegel
called the higher unity, in which oppositions were united
or merged. But whereas Hegel taught that such unities
are formed with logical necessity, Spencer cannot guarantee
such a progress. And no more than other concepts
has the concept of evolution validity beyond experience.

In his ethics \textit{(Data of Ethics.} 1879) Spencer builds (like
Hume and Comte, by the way also Darwin) on sympathy
as foundation and seeks to show how from it there can be grounded
a conduct of life in which human evolution
reaches its most perfect form: as concentration the
ethical life shows itself in self-control, as differentiation in the
rich unfolding of the faculties that compassion with others gives
occasion to, and as determination in conscientiousness
and justice.

\closerule

\chapterheadone{VIII. RECENT TREATMENTS OF PROBLEMS.}{Recent Treatments of Problems}

In romantic idealism and positivist realism two tendencies of
thought stand opposed to each other, which have it in common
that individual problems are for the most part discussed from general
points of view. In a certain contrast to both tendencies
stand endeavors that dwell in detail on individual problems,
perhaps on a single one, and let the general points of view
recede into the background or come as a final
% --- p. 82 ---
\setcounter{footnote}{0}%
\opage{82}outlook. Arranging the exposition according to the characterization of
the four principal problems given in the Introduction, we give
in what follows a survey of such endeavors in the
latter part of the nineteenth century and at the beginning of
the twentieth.

\sectionhead{a. The Psychological Problem.}

α. Of great importance here was the founding of an experimental
psychology by \textsc{Gustav Theodor Fechner} (1801---
1887), who passed from physics to philosophy. In his world-view
he stands close to Romanticism, in that he boldly employs
poetic analogies; but he unites with this idealism
an exact scientific outlook. Since, in his view, the law of the conservation of
energy makes it impossible to assume
a causal relation between the mental and the material, he assumes
that the relation is like that between the inner and the
outer side of a spherical surface: the idealist sees only the former, the materialist
only the latter, and yet they really see the same thing. But
Fechner now sets himself the task of finding out how these
two points of view are related to each other, and it is by this
way that he came to experimental psychology \textit{(Elemente der
Psychophysik.} 1861). Like Kepler, with whom he is
akin in spirit, he worked his way from poetic speculation
to exact investigations. By psychophysics is understood the investigation
of the relation between the mental and the material. By
investigating the relation between the physical stimuli and
the corresponding sense-sensations Fechner arrived at the result
that stimulus and sensation are not directly proportional,
but that the stronger a stimulus is, the more
must it be increased if a certain change of sensation is to arise.
The stimuli must grow in a quotient series in order that
the sensations may grow in a difference series (See Psykologi
V A, 3). By his own experiments and those of others he found this
law confirmed. Yet it is (so far as it holds) hardly, as
Fechner thought, a law for the relation between mental life and
brain life. These are probably (which also agrees
% --- p. 83 ---
\setcounter{footnote}{0}%
\opage{83}best with Fechner's own identity hypothesis) directly proportional,
whereas Fechner's law expresses the relation between
soul or brain on the one side, the physical stimuli
on the other.

What Fechner had founded was carried further by \textsc{Wilhelm} \textsc{Wundt}
(b. 1832), who passed from physiology to philosophy, and
who organized the new science by his work \textit{Physiologische
Psychologie} (1874). A whole series of other problems besides
that treated by Fechner was now taken up, such as the speed
with which psychical phenomena arise, the development
of the perception of space, the laws of the association of representations, etc. --- A
series of investigators have continued what Fechner had founded
and Wundt organized. In Denmark \textsc{Alfred} \textsc{Lehmann}
has worked in this field.

In France the attempt has been made especially to illuminate psychology
by the study of diseased mental life, whose conditions and manifestations often
let the peculiarity of mental life appear in a particularly
striking way. The famous neuropathologist \textsc{Charcot} and
his school gained great influence in this respect; in the
most recent time \textsc{Pierre} \textsc{Janet}, who has passed from philosophy
to psychiatry, works in this direction.

β. Descriptive psychology found eminent practitioners in
\textsc{William} \textsc{James} (b. 1842, d. 1910) \textit{(Principles of Psychology}
1890) and \textsc{Henri} \textsc{Bergson} (b. 1859) (\textit{Les données immédiates de
la concience.} 1888), who both lay great stress on the involuntary
mental life in its continuous stream, in contrast to
the differences and oppositions that practical and theoretical wants
induce us to set up. Already in an essay
of 1884 James had pointed out the original mental
stream of life (the original flux of life, the stream of thought)
as the foundation from which our explicit consciousness continually
draws. Bergson stresses in particular that all quantitative differences
are set up only as means to a clear and distinct apprehension and
understanding of what is immediately given, whose continuous course
presents only qualitative differences. With great art of presentation
the two psychologists attempt to open our eyes to
what is originally and immediately given, which in ordinary
% --- p. 84 ---
\setcounter{footnote}{0}%
\opage{84}language and later in science is parceled out in a mechanical
way.

γ. With regard to the peculiar nature of the life of consciousness
the conception is more and more asserted, which Leibniz and
Kant had expressed, that we must seek this peculiarity
in the fact that every manifestation of mental life makes itself known as a gathering together,
a synthesis, of a more or less manifold content,
whose elements get their character determined by this
very gathering together. This hypothesis about the nature of the life of consciousness
stands in contrast partly to the scholastic doctrine of
a soul-substance (which was upheld by Descartes and his spiritualist
successors), partly to the atomistic conception of
mental life, which was upheld by the English ``association psychology''
(Hume, Stuart Mill). Wilhelm Wundt adheres to it;
it is sought to ground it in the presentation of psychology
which the author of this survey has given, just as it
is also asserted by Pierre Janet on the basis of his
investigations of diseased mental life. (Cf. Psykologi II, 5;
V B, 5). (In his book \textit{La synthèse mentale,} 1908, \textsc{Dwelshauer}
has presented the concept of synthesis as it appears
in the three authors named and in Leibniz and Kant).

The concept of synthesis leads naturally to emphasizing mental life as
a peculiar energy and thereby also to bringing out
the life of will as the chief side of the life of consciousness. On this
point Schopenhauer has exercised a thoroughgoing influence.
A voluntarism is asserted in psychological
theory, in contrast to intellectualism and emotionalism. In
\textsc{Alfred} \textsc{Fouillée}'s \textit{Psychologie des idée-forces} (1893) this
standpoint has found its most interesting expression. All
mental phenomena are symptoms of a need or
striving (appétition), which appears as pleasure or pain,
according as it is favored or hindered. Discriminating and
preferring (discernement et préférence) are originally one,
and every sense perception, every act of thought denotes a more or
less conscious and determinate direction of our impulse of life.

% --- p. 85 ---
\setcounter{footnote}{0}%

\opage{85}\sectionhead{b. The Problem of Existence.}

The philosophy of Romanticism had believed that it could reform natural science.
And yet, precisely at the same time as the origin of this
philosophy, at the end of the eighteenth century,
natural science had made great progress. The old conviction
of the conservation of matter was confirmed by \textsc{Lavoisier}
by weighing, that is, by means of the quantitative method, and a
series of investigators (Priestley, Saussure, and others) demonstrated the main laws
of metabolism in the plant and animal world and thereby incorporated
organic life into the great cycle of nature.

Natural science received a new upswing about the year
1840, notably through \textsc{Robert} \textsc{Mayer}'s discovery of the conservation of
energy (1842). Thoughts which Descartes, Huyghens
and Leibniz had already expressed now received empirical confirmation,
in that it turned out that it is not only in the purely mechanical
domain that no force perishes or arises, but that this
also holds in the interaction among the special forces of nature
themselves, in that a definite quantitative
relation (equivalence) can be established between the manifestation of force (e.g.
motion) that disappears and that (e.g. heat) which arises.

To this was added (1859) Darwin's hypothesis of the
origin of species. Natural science had thus demonstrated, in a magnificent
manner, inner connection where one had previously seen only
fragments and leaps. The question now was what
influence these results would have on the general
world-view. In Herbert Spencer we have already
met an attempt to answer this question. Likewise
in Fechner, who renewed Spinoza's identity hypothesis. In a
brilliant exposition \textsc{Albert Lange} (1828---1875)
\textit{(Geschichte des Materialismus} 1865) sought to demonstrate the impossibility
and the unjustifiability of inserting psychical links into the series
of material causes, whereby the identity hypothesis becomes the
only possibility. Materialism is untenable, because the material
processes from which it would derive the mental phenomena
are themselves known to us only through sensations and
% --- p. 86 ---
\setcounter{footnote}{0}%
\opage{86}representations, that is, through mental processes. Matter is a
concept formed according to the laws that hold for our
sense perception and thinking. Thus Lange combines the problem of existence
and the problem of knowledge.

But besides the Spinozistic hypothesis (which is in the main
upheld by Wundt and is also sought to be grounded in the
presentation of psychology that the author of this survey has
given) there also appear attempts in a materialist and in a spiritualist
direction.

In German literature the upswing of natural science gave occasion
to a series of materialist writings, which had the
merit of spreading scientific knowledge and
ideas in wider circles. With the opposition to dogmatics and
spiritualist speculation there was joined, among the German materialists
of the nineteenth century, as among the French, akin to them in spirit,
in the middle of the eighteenth century,
an idealist interest in humanity and progress.
An idealist disposition is quite compatible with theoretical materialism:
the materialist can acknowledge the high value of mental phenomena
and endeavors, even though he regards
them as mere atomic vibrations. --- The most eminent
of these authors was the physiologist Jacob Moleschott (1822
---1893). Born in Holland, he worked in his youth as
Docent in Heidelberg, and when he was dismissed there on account of
his views, he went to Zürich and later to Italy,
where he carried on significant physiological teaching activity.
In his book \textit{Der Kreislauf des Lebens} (1852) he regards
chemistry as the highest science, because it shows
how matter, and with matter life, with life again thought,
passes through its great cycle. In his autobiography (\textit{An
meine Freunde. Lebenserinnerungen.} 1895) he presents
the history of his thoughts and declares that he has really only
fought against dualism, and that his doctrine, on account of the
inner connection between mind, force and matter, might just as well
be called idealism as materialism. He thus really adheres
to the identity hypothesis. Occasionally the physician \textsc{Louis} \textsc{Büchner} (1824---1899)
also expresses himself in this direction,
% --- p. 87 ---
\setcounter{footnote}{0}%
\opage{87}although he believes he upholds pure materialism. His book
\textit{Kraft und Stoff} (1855) is very unclear as regards the questions
of principle, but is written in a popular language and
won a very large circle of readers.

In contrast to these authors \textsc{Hermann} \textsc{Lotze}
(1817---1881) set himself the task of holding fast the idealism of the Romantic
period in connection with, and on the basis of,
the methods and results of modern natural science. He
began his scientific activity at once as
natural scientist and philosopher, but later passed wholly over to philosophy,
for which he worked a long series of years at the
University of Göttingen. As a natural scientist it was his
endeavor to treat medicine and physiology as mechanical
natural sciences without the then still frequent appeal
to a special ``vital force.'' His philosophical ideas he has
developed in \textit{Mikrokosmus} (1864---68) and in \textit{Drei Bücher
der Metaphysik} (1879).

Whereas the philosophy of Romanticism went the way of deduction and construction,
Lotze is clear that the means of thought which
are at our disposal in grounding our world-view
are analysis and analogy. --- The mechanical view of nature
regards all phenomena as determined by the interaction of
elements which are extended, even though their extension is
extraordinarily small. But what is extended must consist of
parts, cannot be absolutely simple. And extension is moreover
a property that must be explained according to the same principle
as all sense qualities, namely by an interaction of elements;
but the elements whose interaction is to explain
the quality of extension cannot themselves be extended, must
therefore be simpler than the elements at which natural science
stops; they can then only be points of force.
And if they are to stand in interaction with one another, they
cannot be absolutely independent and distinct; they must
be members of a primal substance, which embraces them all and makes
it possible that the state of one can determine that of another. ---
Further than to the indeterminate concepts of unextended elements
and of a primal substance we cannot come by the way of
% --- p. 88 ---
\setcounter{footnote}{0}%
\opage{88}analysis. Only by the way of analogy can we get a positive
determination of these concepts. Lotze agrees with Leibniz
that all metaphysics rests on analogy. Starting from the view that
existence in itself must be conceived either in analogy with the
material or with the mental, he maintains that this
latter analogy is the only possible one, because the mental is
given to us immediately in our own nature, whereas we know only
the material, including our own body and our own brain,
as object. Our body is only the part of the external world
that is nearest to us. The elements of existence and the
primal substance that reveals itself in this interaction
must then be thought of as animated in different degrees. Only in
this way do we get a rational world-view.

The relation between soul and body presents for Lotze
no difficulty, since, according to his metaphysical idealism, they
consist of homogeneous elements. (See further on his standpoint
in this respect Psykologi II 8 c).

\sectionhead{c. The Problem of Knowledge.}

Toward the end of the nineteenth century the insight
gained ground that both the Romantic and the positivist
reaction against Kant's critical philosophy had gone too far.
The critical tendency had not disappeared in the course of the century,
but now became of greater importance, and one immersed
oneself again with zeal in Kant's writings. In Germany \textsc{Albert} \textsc{Lange}
in particular paved the way for the so-called neo-criticism or
neo-Kantianism. In France \textsc{Charles} \textsc{Renouvier}
(1815---1903) stands as the energetic renewer of the critical
philosophy.

Kant's task had been to find the fundamental concepts and
the presuppositions on which a scientific understanding
of existence rests. In a certain sense one can say that
these fundamental concepts and presuppositions are the means by
whose help we can attain the goal that our intellectual interest
lets us strive after. But the fundamental concepts (space, time, magnitude,
causality) were for Kant at the same time expressions of the
% --- p. 89 ---
\setcounter{footnote}{0}%
\opage{89}human mind's forms, of our intellectual nature, and it
was for him no accident that in science we work precisely
with these fundamental concepts. In this \textit{criticism} differs
from an epistemological conception that has appeared in
recent years, and which grounds the fundamental
concepts and presuppositions of science exclusively in the utility, the theoretical
and practical significance, that they have. This conception, the
so-called \textit{pragmatism,} is maintained by William James in close
connection with his psychological conception mentioned above,
according to which we leave the involuntary and immediate flow
in which our mental life manifests itself only when necessity
requires it. The thoughts we develop and the
propositions we set up are always called forth by definite interests.
Were we wholly indifferent, we would not form definite
concepts or draw definite conclusions. (W. \textsc{James}:
\textit{Pragmatism.} 1907).

A pragmatic conception was also asserted from the
side of natural science. Eminent natural scientists, who reflected
on the principles of their science, became aware
that the determinations of concepts and the presuppositions of science must
seek their justification in this, that they make possible an intellectual
elaboration and presentation of the facts. Only
on this rests their necessity, which is therefore not unconditional,
so long as the possibility is not excluded that other
concepts and presuppositions than those now employed might just as
well lead to the same goal. The famous English physicist
Maxwell expressed himself in this direction as early as 1855, and
later \textsc{Ernst} \textsc{Mach}, on the basis of studies in the history of
physics, asserted the same point of view. Both
investigators lay great stress on the analogies by whose help
science transfers representations from one domain of experience
to another; thus the mechanical view of nature rests
on a great, carried-through analogy between the motion of masses
in space and the qualitative changes of things
(temperature, electricity, etc.) (Ernst Mach: \textit{Die Mechanik
in ihrer} Entwickelung. 1883. --- \textit{Erkenntnis und} Irrtum. 1905).
Mach uses, instead of the word pragmatism, rather the expression
% --- p. 90 ---
\setcounter{footnote}{0}%
\opage{90}economical epistemology, laying stress on the fact that we
restrict our definitions and hypotheses to what is strictly
necessary.

\textsc{Richard} \textsc{Avenarius} (1843---1896) has presented the
economical or pragmatic theory of knowledge from a psychological
point of view, in that he has sought to give what one might
call the natural history of problems. \textit{(Philosophie als} Denken
\textit{der Welt nach dem Prinzip des kleinsten Kraftmaasses.} 1876.
\textit{Kritik der reinen} Erfahrung. 1888---1890). A problem presupposes
a tension between the individual and the surrounding world.
Such a tension arises when the stimuli that
come from the external world demand either a greater or a
smaller discharge of energy than the individual is able to supply.
When, on the other hand, stimulus (P) and energy (E) correspond to each other
(so that P = E), there is no problem and no unrest;
the individual recognizes its external world, feels at home, and has
confidence in its representations and observations.

If a greater labor is demanded than the individual can supply
(that is, P > E), the individual feels a stranger in its world.
Recognition is not possible, and the given stands as full of
contradictions, deviations and exceptions. With every extension
of experience and of the mental horizon, new
problems therefore arise. Rising culture brings an increase of problems.
--- --- If, conversely, more energy is at hand than the external world
demands (so that P < E), there will arise a need
that craves beyond the given, a romantic ferment or a
practical idealism will appear. The unrest will persist until
either the energy is exhausted or the external world is remolded
according to the need of the individual.

The equalization (deproblematization) is not always a real
solution. Perhaps only a provisional or purely
individual point of view is reached, without the possibility of a new tension
(problematization) being excluded. Definitive and universal
deproblematization comes about only when a complete adaptation
to the given is attained, and all provisional and individual
elements in the conception are eliminated. This is attained
when knowledge essentially consists in quantitative description,
% --- p. 91 ---
\setcounter{footnote}{0}%
\opage{91}which constantly shows the consequent to be the equivalent of the
antecedent. Then the standpoint of pure experience is reached. ---

In Danish literature K. \textsc{Kroman}'s book \textit{Vor Naturerkendelse}
(1883) points in a pragmatic direction, in so far as it finds in
the impulse of self-preservation the psychological foundation for
the postulates that set up the principles of scientific
knowledge (the principle of identity, the principle of causality). ---

Whereas criticism teaches that we set up the fundamental concepts
and presuppositions \textit{that, according to the nature of our mind, we must set up,}
pragmatism teaches that we set up the fundamental concepts and presuppositions
\textit{that we have need of.} It is clear that the means we
have at our disposal for solving our tasks must be determined by
our nature and depend on it; there must be an inner
connection between the means of our knowledge and the
nature of our mind. Synthesis is at once the means to knowledge and
the form of the involuntary working of our mind. Pragmatism adheres
to the occasions that lead to the development of our conditions
(our ``forms''); criticism is inclined to regard
the ``forms'' as given once for all.

\medskip\noindent Appendix. \textbf{The History of the Causal Principle in Modern Times.}

1. The principle is set up as self-evident, even in the definite
form that there cannot be More in the effect than in
the cause. No distinction is made between ground (ratio) and cause
(causa). (\textsc{Descartes} and \textsc{Spinoza}).

2. \textsc{Hume}'s criticism of the causal principle: this can be proved neither
by inference nor by observation, but rests
only on custom or acquired instinct.

3. \textsc{Kant}: The concept of cause is a category, one of the fundamental forms
of our thinking, a special kind of synthesis. The causal principle
is proved by the fact that only by means of it can we distinguish
between mere representation and real experience; if there is to be
experience in the strict sense (as distinct from mere representation),
the concept of cause must hold. There is only analogy between
ground and cause.

4. The concept of cause is a means to knowledge and
% --- p. 92 ---
\setcounter{footnote}{0}%
\opage{92}orientation in the world; the causal principle is merely a working hypothesis.
\textit{(Pragmatism).}

5. \textit{The conception presented in Psykologi V D, 3} agrees
with Kant that the concept of cause is one of the forms in which
synthesis, the formal unity of consciousness, displays itself,
--- and agrees with Hume that no conclusive
proof can be given of the validity of the causal principle. Over against pragmatism
it is stressed that the means and ways our knowledge
employs may very well be expressions of the nature of our mind,
indeed, must be so, even though the forms of our nature cannot be
found and fixed once for all. In the activity of thought
an essential side of our nature expresses itself, and such an
expression is both end and means. The intellectual interest
has its independent significance in the human
life of the spirit.

\sectionhead{d. The Problem of Valuation.}

α. In the domain of ethics the conflict toward the end
of the nineteenth century was between Kant's morality of reason
and the morality of utility maintained by Bentham and Stuart Mill;
the former held to the inner, free submission to
a universally valid law, the latter to the demonstrable
effects of actions. In an acute book \textit{(Methods of Ethics} 1877)
\textsc{Henry} \textsc{Sidgwick} (1838---1900) demonstrated that within
the morality of utility a sharp distinction had to be made between two standpoints:
the one takes account only of the individual's
own utility, the other understands by utility the universal
utility, the utility of the whole society. What is useful from the
one standpoint is not always useful from the other standpoint,
and a difficult, perhaps insoluble problem can arise here.
In a masterly analysis of generally recognized
practical morality (the morality of common sense) Sidgwick shows
that it has a universal morality of utility as its presupposition.

In the presentation of ethics which the author of this survey
published in 1887, it was shown that Kant's ethics and Bentham's and Mill's
morality of utility, in spite of all opposition, rested on a common
% --- p. 93 ---
\setcounter{footnote}{0}%
\opage{93}presupposition, namely that human actions are to be valued
according to a universally valid standard. This very standard (the
fundamental value that determines the individual values) can of course
not in turn be valued. It is the business of historical and psychological
investigation to show how a standpoint
that recognizes and applies such a standard comes into being;
it is the business of ethics to apply the presupposed standard to
the particular circumstances of life. In the very presupposition of a universally valid
standard it lies that individual ethics must be
a part of social ethics. The great problem is how a single
individual can be led over from one fundamental value to another.

In the most recent time the social character of ethical
valuation has been strongly impressed especially by French investigators.
Already in Comte ethics was a part of sociology; only
in his later period did he separate it off. The French sociological
school with \textsc{Émile} \textsc{Durkheim} at its head lays
social solidarity, grounded partly on the bond of blood, partly on
community of labor, at the foundation of morality. The individual's conscience
is determined, consciously and unconsciously, by the influence of society
through traditions, manners, institutions and discipline.
(\textsc{Durkheim}: \textit{La Division du Travail Social.} 1893). --- Whereas
Durkheim predominantly investigates the historical
and psychological development of practical morality, \textsc{Léon} \textsc{Bourgeois} and
his followers proceed from the view that every individual, by virtue of the solidarity
of society, stands in debt to the earlier generations,
a debt he must repay by working on the tasks of his own time.
From solidarity spring society's right and the individual's duty.
Léon Bourgeois's book \textit{Solidarité} has appeared in Danish translation
(1910). --- In Danish literature N. \textsc{Bang}'s line of thought
in \textit{Begrebet Moral} (1897) goes in a similar direction.

β. As Avenarius demonstrated, practical problems arise
when there stirs in the individual a greater fullness of
energy than the surroundings can lay claim to. A similar
thought stirs in various authors of the most recent time,
who strongly emphasize spontaneity, the involuntary striving and
need for expansion. In the descriptive psychology of James and Bergson
this side of mental life comes strongly forward, and it
% --- p. 94 ---
\setcounter{footnote}{0}%
\opage{94}has also been applied to the problem of valuation, most passionately
and ruthlessly by \textsc{Friedrich} \textsc{Nietzsche} (1844---
1900), who demanded an entirely new valuation of life, with a revaluation
of all values hitherto recognized. In his early work \textit{Die
Geburt der Tragœdie} (1872) he sets the tragic-poetic
conception of life, represented by Dionysus and Apollo, in opposition
to the intellectual optimism of Socrates. Science
(Socrates) stands below art (Apollo), which in turn stands below
overflowing life (Dionysus), which even takes up the struggle
with suffering and death. But the vigorous unfolding of life,
without which existence does not acquire positive value, he found
possible only when a ruling estate represents strength
and self-assertion, while a circle of workers and philistines
(die Heerde) attends to the subordinate labor.
Aristocracy exists for its own sake, since it expresses
the highest strength of life and immediate value; in it the goal of history
is attained, while ``the herd'' is only copy, means or hindrance
\textit{(Jenseits von Gut und} Böse. 1886. --- \textit{Die Genealogie
der Moral} 1887). In the poetic work \textit{Also sprach Zarathustra,}
especially in the sketch for its conclusion, the possibility is nevertheless hinted
that this social dualism might be overcome when
the revaluation was completed.

Related thoughts have also been expressed from other standpoints.
--- \textsc{Jean} \textsc{Marie} Guyau's (1854---1888) main thought
was life's need of unfolding, communication and devotion,
the spontaneous working of the faculties of the soul, which ever anew leads
beyond results attained and restraining barriers. In this
need Guyau finds the possibility of sympathy (altruism) and
magnanimity, which the English school explained by association of representations
or by the struggle for life. \textit{(Esquisse d'une Morale
sans obligation ni sanction.} 1885). And out of this need
there may also develop an exalted and courageous mood of life,
which dares to take up the struggle with the disharmonies of existence,
even though it cannot lean on supernatural guarantees.
(\textit{L'irréligion de l'avenir.} 1887). --- The author of this
historical survey has sought in his Religionsfilosofi (1901) to
show that behind all magic and religion lies a need to
% --- p. 95 ---
\setcounter{footnote}{0}%
\opage{95}maintain the constant subsistence of those goods (values) which
for man at a given stage are the highest. This need does not have
to fall away, even though it can no longer be given shape
in myth, dogma and cult, but must seek expression in
poetic symbols, and even though reflection leads to the insight
to which the history of values under the development of the human species
already leads, that our values cannot have lasting
subsistence in the form we know. There can be a
continuous development in the domain of values, as in other
domains, and it is this that can be an object of our
faith. --- --- \textsc{William James} shows in his work \textit{Varieties of religious
experience} (1902), supported by a study of religious phenomena
as well as by what he had demonstrated in his psychology,
how small a part of our mental being can be explained
clearly and definitely. Through many degrees the conscious goes
back into the unconscious, or the subconscious, and our first
presuppositions often stem from this ``subliminal'' region,
from what lies beneath the threshold of consciousness. An involuntary
conviction is the deepest thing in us (the deep thing in us).
James is inclined to regard the influences from that dim
region as the ways by which a higher order of things
works in us. Every closer determination of the origin of these influences
is an interpretation that will turn out differently
in different individuals. Most people are too inclined
to a dogmatic formulation of their involuntary need.
Religion is, according to James, man's total mood toward
life, a mood that contains in itself the sum of man's
experiences of life. Thereby a unity and harmony is attained
in the mental life, which would otherwise be lacking. And the fullness of life
that can stream from the unconscious sources and continually
nourishes our conscious life can be strong enough that tragedies
and shipwrecks never for us denote the last word of experience.

\closerule

\end{document}
